\chapter{پردازش متن}

تمامی سیستم‌عامل‌های شبه‌یونیکس (\en{Unix-like}) برای ذخیره‌سازی پیکربندی‌ها و داده‌ها، وابستگی شدیدی به فایل‌های متنی ساده (\en{Plain Text}) دارند. از این رو، وجود مجموعه‌ی گسترده‌ای از ابزارها برای مدیریت و پردازش متن در لینوکس کاملاً منطقی است. در این فصل، بر برنامه‌هایی تمرکز می‌کنیم که جهت «قطعه‌بندی و بازآرایی» (\en{Slicing and Dicing}) محتوای متنی به کار می‌روند. در فصل آینده، جنبه‌های پیشرفته‌تری از پردازش متن را بررسی خواهیم کرد که عمدتاً بر قالب‌بندی (\en{Formatting}) متون جهت چاپ یا ارائه در فرمت‌های خوانا برای انسان متمرکز هستند.

در این فصل، با ترکیبی از ابزارهای کلاسیک و مدرن آشنا خواهیم شد:

\begin{itemize}
  \item \cmd{cat}: الحاق فایل‌ها و نمایش محتوای آن‌ها در خروجی استاندارد (\en{Standard Output}).
  \item \cmd{sort}: مرتب‌سازی خطوط فایل‌های متنی.
  \item \cmd{uniq}: گزارش یا حذف خطوط تکراریِ مجاور در یک فایل.
  \item \cmd{cut}: استخراج بخش‌های مشخصی از هر خط یک فایل.
  \item \cmd{paste}: ادغام خطوط چندین فایل به‌صورت افقی.
  \item \cmd{join}: پیوند دادن خطوط دو فایل بر اساس یک فیلد (زمینه) مشترک.
  \item \cmd{tac}: الحاق و نمایش فایل‌ها از انتها به ابتدا (معکوسِ \cmd{cat}).
  \item \cmd{rev}: معکوس کردن ترتیب کاراکترهای هر خط.
  \item \cmd{comm}: مقایسه‌ی خط‌به‌خط دو فایل مرتب‌شده.
  \item \cmd{diff}: مقایسه‌ی خط‌به‌خط دو فایل و گزارش تفاوت‌های آن‌ها.
  \item \cmd{patch}: اعمال تغییرات تولیدشده توسط \cmd{diff} بر روی فایل اصلی.
  \item \cmd{tr}: جایگزینی (ترجمه) یا حذف کاراکترها.
  \item \cmd{sed}: ویرایشگر جریانی (\en{Stream Editor}) برای پالایش و تغییر متون.
  \item \cmd{aspell}: بررسی‌کننده‌ی تعاملی غلط‌های املایی.
\end{itemize}

\section{کاربردهای متن در لینوکس}

تا به این مرحله، با ویرایشگرهای متنی نظیر \cmd{nano} و \cmd{vim} آشنا شده‌ایم، فایل‌های پیکربندی متعددی را بررسی کرده و خروجی دستورات فراوانی را مشاهده کرده‌ایم که همگی در قالب متن ارائه شده‌اند. اما قلمرو کاربرد متن فراتر از این موارد است. در واقع، متن ساده ستون فقراتِ دنیای رایانش محسوب می‌شود.

\subsection{مستندات}

بسیاری از متخصصان، اسناد خود را با بهره‌گیری از قالب‌های متنی ساده (\en{Plain Text}) تدوین می‌کنند. اگرچه کاربرد یک فایل متنی کوچک برای یادداشت‌های روزمره بدیهی است، اما نگارش اسناد حجیم و پیچیده نیز در این قالب کاملاً رایج است. یکی از راهبردهای استاندارد، نگارش محتوا در قالب متن ساده و سپس استفاده از زبان‌های نشانه‌گذاری (\en{Markup Languages}) برای تعیین ساختار و قالب‌بندی نهایی است.

بسیاری از مقالات علمی و متون فنی با این روش به رشته‌ی تحریر درمی‌آیند؛ چراکه سیستم‌های پردازش متن در یونیکس جزو نخستین بسترهایی بودند که امکانات پیشرفته‌ی تایپوگرافی را در اختیار نویسندگان قرار دادند. در سال‌های اخیر، زبان نشانه‌گذاری \en{Markdown} به دلیل سادگی و کارایی بالا در تبدیل متن به اسناد خوش‌ساخت، محبوبیت چشمگیری یافته است.

\subsection{صفحات وب}

شاید بتوان گفت صفحات وب، پرکاربردترین نوع اسناد الکترونیکی در جهان امروز هستند. این صفحات در واقع اسنادی متنی هستند که با استفاده از زبان‌های نشانه‌گذاری همچون \en{HTML} (\en{Hypertext Markup Language}) یا \en{XML} (\en{Extensible Markup Language})، ساختار بصری و چیدمان محتوا را برای مرورگرها توصیف می‌کنند.

\subsection{پست الکترونیک}

پست الکترونیک در ذات خود یک رسانه‌ی مبتنی بر متن است. حتی پیوست‌های غیرمتنی نیز جهت انتقال در شبکه، به کدهای متنی تبدیل می‌شوند. با دانلود یک پیام ایمیل و باز کردن آن در ابزاری نظیر \cmd{less} به‌وضوح می‌توان این ساختار را مشاهده کرد. هر پیام با یک «سرآیند» (\en{Header}) آغاز می‌شود که مشخصات فرستنده و مسیرهای پیموده‌شده توسط پیام را توصیف می‌کند و در ادامه، «بدنه‌ی پیام» (\en{Body}) حاوی محتوای اصلی قرار می‌گیرد.

\subsection{خروجی چاپ}

در سیستم‌های شبه‌یونیکس، داده‌های ارسالی به چاپگر معمولاً به‌صورت متن ساده (\en{Plain Text}) یا با استفاده از یک زبان توصیف صفحه (\en{Page Description Language}) با ساختار متنی منتقل می‌شوند. در گذشته، \en{PostScript} زبان استاندارد و محوری این فرایند بود؛ اما امروزه \en{PDF} به‌عنوان قالب اصلی و ترجیحی برای ارسال داده‌های چاپی جایگزین آن شده است. در نهایت، این کدهای متنی (چه \en{PostScript} و چه \en{PDF}) به مفسری مانند \en{Ghostscript} فرستاده می‌شوند تا داده‌های گرافیکی (\en{Raster}) موردنیاز برای چاپ نهایی را تولید کنند.

\subsection{کد منبع نرم‌افزار}

بسیاری از ابزارهای خط فرمان در محیط‌های شبه‌یونیکس با هدف تسهیل مدیریت سیستم و توسعه‌ی نرم‌افزار طراحی شده‌اند؛ ابزارهای پردازش متن نیز از این قاعده مستثنی نیستند. علت کلیدی اهمیت پردازش متن برای توسعه‌دهندگان این است که هر نرم‌افزاری در نخستین گامِ حیات خود، یک متن ساده است. کد منبع (\en{Source Code}) که مستقیماً توسط برنامه‌نویس نگاشته می‌شود، همواره در قالب متن ذخیره و مدیریت می‌گردد.

\section{بازبینی ابزارهای پایه}

در فصل ششم، با برخی دستوراتی آشنا شدیم که علاوه بر آرگومان‌های خط فرمان، توانایی دریافت داده از «ورودی استاندارد» (\en{stdin}) را نیز دارند. در آن بخش صرفاً به معرفی اجمالی آن‌ها بسنده کردیم، اما اکنون به بررسی دقیق‌تر نحوه‌ی به‌کارگیری این ابزارها در فرآیندهای پردازش متن خواهیم پرداخت.

\subsection{دستور cat: مدیریت نویسه‌های کنترلی}

برنامه \cmd{cat} دارای گزینه‌های کاربردی متعددی است که فراتر از نمایش ساده‌ی متن، به بازبینی دقیق محتوا کمک می‌کنند. یکی از مهم‌ترین این گزینه‌ها، گزینه \cmd{-A} است که برای آشکارسازی نویسه‌های غیرقابل‌چاپ (\en{Non-printing characters}) به‌کار می‌رود. در بسیاری از موارد، شناسایی نویسه‌های کنترلی پنهان، مانند نویسه‌ی \en{Tab} (که با فاصله متفاوت است) یا نویسه‌ی بازگشت به ابتدای سطر (\en{Carriage Return} که در فایل‌های متنی سیستم‌عامل \en{MS-DOS} رایج است)، اهمیت حیاتی دارد. همچنین این ابزار برای تشخیص فاصله‌های ناخواسته‌ی انتهای خط (\en{Trailing Spaces}) بسیار مفید است.

برای درک بهتر، بیایید از \cmd{cat} به‌عنوان یک واژه‌پرداز ابتدایی استفاده کنیم. با هدایت خروجی \cmd{cat} به یک فایل، متن دلخواه را تایپ کرده، در ابتدای سطر یک کلید \en{Tab} فشار می‌دهیم، در انتها چند فاصله‌ی اضافی وارد کرده و در نهایت با فشردن \cmd{Ctrl+d} (نشان‌دهنده‌ی پایان فایل یا \en{EOF})، عملیات را خاتمه می‌دهیم:

\begin{latin}
\begin{lstlisting}
[me@linuxbox ~]$ cat > foo.txt
   The quick brown fox jumps over the lazy dog.
[me@linuxbox ~]$
\end{lstlisting}
\end{latin}

حال با استفاده از گزینه‌ی \cmd{-A} محتوای فایل را بازبینی می‌کنیم:

\begin{latin}
\begin{lstlisting}
[me@linuxbox ~]$ cat -A foo.txt
^IThe quick brown fox jumps over the lazy dog.  $
[me@linuxbox ~]$
\end{lstlisting}
\end{latin}

نتایج فوق نشان‌دهنده‌ی دو نکته‌ی فنی مهم است: نخست، نویسه‌ی \en{Tab} با نماد \cmd{\^{}I} نمایش داده شده است (این یک قرارداد متداول در یونیکس برای نمایش \cmd{Ctrl+i} یا همان نویسه‌ی تب است). دوم، نماد \cmd{\$} در انتهای سطر ظاهر شده که دقیقاً مرز پایانی خط را نشان می‌دهد؛ وجود فاصله میان نقطه‌ی انتهایی جمله و نماد \cmd{\$} مؤید وجود فاصله‌های اضافی در انتهای خط است.

\section{تمایز متون MS-DOS و Unix}

یکی از دلایل کلیدی استفاده از ابزار \cmd{cat} برای پایش نویسه‌های غیرقابل‌چاپ، شناسایی کاراکترهای پنهانی است که ممکن است در فایل‌های متنی ما پنهان شده باشند. یکی از رایج‌ترین نمونه‌های این نویسه‌ها، نویسه‌ی «بازگشت به ابتدای سطر» (\en{Carriage Return}) است که در فایل‌های سیستم‌های \en{MS-DOS} و ویندوز یافت می‌شود. علت این امر، تفاوت میان استانداردهای پایان خط در یونیکس و \en{DOS} است. در سیستم‌های یونیکس، پایان هر سطر با نویسه‌ی خط جدید (\en{Linefeed} – کد \en{ASCII 10}) مشخص می‌شود. در حالی‌که در \en{MS-DOS} و سیستم‌های جانشین آن، پایان سطر با توالی دو نویسه‌ی «بازگشت به ابتدای سطر» (\en{Carriage Return} – کد \en{ASCII 13}) و به‌دنبال آن «خط جدید» (\en{Linefeed}) نشان داده می‌شود.

روش‌های متعددی برای تبدیل فایل‌ها از فرمت \en{DOS} به یونیکس وجود دارد. در بسیاری از توزیع‌های لینوکس، ابزارهای تخصصی \cmd{dos2unix} و \cmd{unix2dos} برای این منظور تعبیه شده‌اند. با این حال، در صورت عدم دسترسی به این برنامه‌ها نیز جای نگرانی نیست؛ چرا که فرآیند تبدیل یک فایل از فرمت \en{DOS} به یونیکس، صرفاً شامل حذف نویسه‌های ناخواستهٔ «بازگشت به ابتدای سطر» است که با بهره‌گیری از ابزارهای پردازش متنی که در ادامهٔ این فصل بررسی خواهیم کرد، به سادگی قابل انجام است.

برنامه‌ی \cmd{cat} علاوه بر نمایش، گزینه‌هایی برای اعمال اصلاحات جزئی در خروجی متن ارائه می‌دهد. دو گزینه کاربردی در این زمینه عبارتند از: \cmd{-n} برای شماره‌گذاری سطور و \cmd{-s} (مخفف \en{squeeze-blank}) جهت ادغام و حذف خطوط خالیِ متوالی. مثال زیر نحوه‌ی عملکرد این دو گزینه را نشان می‌دهد:

\begin{latin}
\begin{lstlisting}
[me@linuxbox ~]$ cat > foo.txt
The quick brown fox


jumps over the lazy dog.
[me@linuxbox ~]$ cat -ns foo.txt
     1  The quick brown fox
     2
     3  jumps over the lazy dog.
[me@linuxbox ~]$
\end{lstlisting}
\end{latin}

در این نمونه، فایل \file{foo.txt} به‌گونه‌ای بازنویسی شد که شامل دو خط متن با فاصله‌ی دو خط خالی میان آن‌ها باشد. با اجرای دستور \cmd{cat -ns} مشاهده می‌کنید که خطوط خالیِ مکرر فشرده شده و تمامی سطور خروجی شماره‌گذاری شده‌اند. اگرچه این یک عملیات پردازشی ابتدایی است، اما گویای توانمندی این ابزار در مدیریت جریان‌های متنی است.

\section{دستور sort}

برنامه‌ی \cmd{sort} محتویات ورودی استاندارد (\en{stdin}) یا یک یا چند فایل تعیین‌شده در خط فرمان را مرتب‌سازی کرده و نتیجه را در خروجی استاندارد (\en{stdout}) چاپ می‌کند. با بهره‌گیری از همان تکنیکِ ایجاد فایل که در بخش \cmd{cat} آموختیم، می‌توانیم پردازش مستقیم ورودی از طریق صفحه‌کلید را به شکل زیر مشاهده کنیم:

\begin{latin}
\begin{lstlisting}
[me@linuxbox ~]$ sort > foo.txt
c
b
a
[me@linuxbox ~]$ cat foo.txt
a
b
c
\end{lstlisting}
\end{latin}

در این نمونه، پس از اجرای دستور، حروف \cmd{b}، \cmd{c} و \cmd{a} را وارد کرده و با فشردن کلید ترکیبی \cmd{Ctrl+d} وضعیت پایان فایل (\en{EOF}) را اعلام کردیم. پس از بازبینی فایل \file{foo.txt} مشاهده می‌شود که سطرها بر اساس ترتیب الفبایی مرتب شده‌اند.

از آنجا که ابزار \cmd{sort} قابلیت پذیرش چندین فایل به‌صورت هم‌زمان را دارد، می‌توان از آن برای ادغام و مرتب‌سازی یکپارچه‌ی چندین منبع داده استفاده کرد. برای مثال، جهت ترکیب محتوای سه فایل متنی مجزا در قالب یک فهرست واحد و مرتب‌شده، دستور زیر به‌کار می‌رود:

\begin{latin}
\begin{lstlisting}
sort file1.txt file2.txt file3.txt > final_sorted_list.txt
\end{lstlisting}
\end{latin}

این دستور ابتدا محتوای تمامی فایل‌ها را تجمیع نموده و سپس عملیات مرتب‌سازی را روی کل داده‌ها اعمال می‌کند.

برنامه \cmd{sort} از گزینه‌های متعددی برای کنترل دقیق فرآیند مرتب‌سازی بهره می‌برد. جدول زیر، کاربردی‌ترین گزینه‌های این دستور را معرفی می‌کند:

\begin{table}[htbp]
\centering
\caption{گزینه‌های متداول دستور \cmd{sort}}
\label{tab:sort-options}
\begin{tabularx}{\textwidth}{l l X}
\toprule
\textbf{گزینه کوتاه} & \textbf{گزینه بلند} & \textbf{شرح عملکرد} \\
\midrule
\cmd{-b} & \cmd{--ignore-leading-blanks} & نادیده گرفتن فواصل ابتدای سطر؛ مرتب‌سازی بر اساس نخستین کاراکتر غیرفضای‌خالی انجام می‌شود. \\
\addlinespace
\cmd{-f} & \cmd{--ignore-case} & غیرفعال کردن حساسیت به حروف بزرگ و کوچک (\en{Case-insensitive}). \\
\addlinespace
\cmd{-n} & \cmd{--numeric-sort} & مرتب‌سازی خطوط بر اساس ارزش عددیِ رشته‌ها، به‌جای ترتیب واژه‌نامه‌ای (\en{Lexicographic}) پیش‌فرض. در حالت پیش‌فرض، دستور \cmd{sort} رشته‌ها را کاراکتر به کاراکتر و بر اساس جدول کدهای \en{ASCII} مقایسه می‌کند؛ در نتیجه، عددی مانند «\en{100}» پیش از عدد «\en{20}» قرار می‌گیرد، چرا که کاراکتر «\en{1}» از نظر الفبایی مقدم بر «\en{2}» است. با فعال کردن گزینه‌ی \cmd{-n}، برنامه‌ی \cmd{sort} محتوای ابتدای هر خط را به‌عنوان یک عدد صحیح یا اعشاری تفسیر کرده و بر اساس مقدار واقعی آن مرتب‌سازی را انجام می‌دهد؛ این ویژگی به‌ویژه برای مرتب‌سازی صحیح فایل‌های حاوی داده‌های عددی (مانند اندازه‌ی فایل‌ها یا شماره‌های شناسایی) ضروری است. \\
\addlinespace
\cmd{-r} & \cmd{--reverse} & معکوس کردن جهت مرتب‌سازی؛ نمایش نتایج به صورت نزولی. \\
\addlinespace
\cmd{-k} & \cmd{--key=field1[,field2]} & تعیین یک کلید خاص برای مرتب‌سازی بر اساس محدوده‌ای از فیلد ۱ تا فیلد ۲، به‌جای کل سطر. \\
\addlinespace
\cmd{-m} & \cmd{--merge} & ادغام فایل‌هایی که از پیش مرتب شده‌اند؛ این گزینه چندین فایل را بدون پردازش مجدد با هم ترکیب می‌کند. \\
\addlinespace
\cmd{-o} & \cmd{--output=file} & هدایت مستقیم خروجی مرتب‌شده به یک فایل مشخص، به‌جای استفاده از خروجی استاندارد. \\
\addlinespace
\cmd{-t} & \cmd{--field-separator=char} & تعیین کاراکتر جداکننده‌ی فیلدها (مانند کاما یا دو‌نقطه). به‌صورت پیش‌فرض از فاصله (\en{Space}) یا تب (\en{Tab}) استفاده می‌شود. \\
\bottomrule
\end{tabularx}
\end{table}

اگرچه اکثر این گزینه‌ها عملکردی بدیهی دارند، اما برخی از آن‌ها نیازمند واکاوی دقیق‌تر هستند. در گام نخست، به بررسی گزینه‌ی \cmd{-n} می‌پردازیم که برای مرتب‌سازی عددی (\en{Numeric Sort}) طراحی شده است. این گزینه به ابزار \cmd{sort} اجازه می‌دهد تا برخلاف حالت پیش‌فرضِ الفبایی، مقادیر را بر اساس وزن عددی آن‌ها پردازش کند.

یک سناریوی کاربردی برای درک این تفاوت، استفاده از دستور \cmd{du} جهت شناسایی دایرکتوری‌هایی است که بیشترین فضای دیسک را اشغال کرده‌اند. به‌صورت عادی، دستور \cmd{du} خروجی خود را بر اساس ترتیب حروف الفبای مسیرها ارائه می‌دهد:

\begin{latin}
\begin{lstlisting}
[me@linuxbox ~]$ du -s /usr/share/* | head
252	/usr/share/aclocal
96	/usr/share/acpi-support
8	/usr/share/adduser
196	/usr/share/alacarte
344	/usr/share/alsa
8	/usr/share/alsa-base
12488	/usr/share/anthy
8	/usr/share/apmd
21440	/usr/share/app-install
48	/usr/share/application-registry
\end{lstlisting}
\end{latin}

در مثال بالا، خروجی را به دستور \cmd{head} پایپ (\en{Pipe}) کردیم تا صرفاً ۱۰ سطر نخست نمایش داده شود. حال برای استخراج فهرستی از ۱۰ اشغال‌کننده‌ی اصلی فضای ذخیره‌سازی، خروجی را با استفاده از ترکیب گزینه‌های \cmd{-nr} (مرتب‌سازی عددی و معکوس) پردازش می‌کنیم:

\begin{latin}
\begin{lstlisting}
[me@linuxbox ~]$ du -s /usr/share/* | sort -nr | head
509940	/usr/share/locale-langpack
242660	/usr/share/doc
197560	/usr/share/fonts
179144	/usr/share/gnome
146764	/usr/share/myspell
144304	/usr/share/gimp
135880	/usr/share/dict
76508	      /usr/share/icons
68072	      /usr/share/apps
62844       /usr/share/foomatic
\end{lstlisting}
\end{latin}

استفاده از ترکیب گزینه‌های \cmd{-n} و \cmd{-r} جهت انجام مرتب‌سازی عددی معکوس، باعث می‌شود بزرگ‌ترین مقادیر در ابتدای خروجی قرار گیرند. این روش زمانی که عدد مورد نظر در ابتدای سطر باشد به‌خوبی عمل می‌کند؛ اما پرسش اینجاست که اگر بخواهیم مرتب‌سازی را بر اساس مقداری در میانه‌ی سطر انجام دهیم، چه راهکاری وجود دارد؟ برای نمونه، خروجی دستور \cmd{ls -l} را در نظر بگیرید:

\begin{latin}
\begin{lstlisting}
[me@linuxbox ~]$ ls -l /usr/bin | head
total 152948
-rwxr-xr-x 1 root root    34824 2016-04-04 02:42 [
-rwxr-xr-x 1 root root   101556 2007-11-27 06:08 a2p
-rwxr-xr-x 1 root root    13036 2016-02-27 08:22 aconnect
-rwxr-xr-x 1 root root    10552 2007-08-15 10:34 acpi
-rwxr-xr-x 1 root root     3800 2016-04-14 03:51 acpi_fakekey
-rwxr-xr-x 1 root root     7536 2016-04-19 00:19 acpi_listen
-rwxr-xr-x 1 root root     3576 2016-04-29 07:57 addpart
-rwxr-xr-x 1 root root    20808 2016-01-03 18:02 addr2line
-rwxr-xr-x 1 root root   489704 2016-10-09 17:02 adept_batch
\end{lstlisting}
\end{latin}

اگرچه خودِ دستور \cmd{ls} قابلیت مرتب‌سازی بر اساس حجم را دارد، اما می‌توانیم این وظیفه را به ابزار \cmd{sort} بسپاریم تا کنترل بیشتری بر فرآیند داشته باشیم. برای مرتب‌سازی این فهرست بر اساس اندازه‌ی فایل، از دستور زیر استفاده می‌کنیم:

\begin{latin}
\begin{lstlisting}
[me@linuxbox ~]$ ls -l /usr/bin | sort -nrk 5 | head
-rwxr-xr-x 1 root root 8234216 2016-04-07 17:42 inkscape
-rwxr-xr-x 1 root root 8222692 2016-04-07 17:42 inkview
-rwxr-xr-x 1 root root 3746508 2016-03-07 23:45 gimp-2.4
-rwxr-xr-x 1 root root 3654020 2016-08-26 16:16 quanta
-rwxr-xr-x 1 root root 2928760 2016-09-10 14:31 gdhtui
-rwxr-xr-x 1 root root 2928756 2016-09-10 14:31 gdb
-rwxr-xr-x 1 root root 2602236 2016-10-10 12:56 net
-rwxr-xr-x 1 root root 2304684 2016-10-10 12:56 rpcclient
-rwxr-xr-x 1 root root 2241832 2016-04-04 05:56 aptitude
-rwxr-xr-x 1 root root 2202476 2016-10-10 12:56 smbcacls
\end{lstlisting}
\end{latin}

اغلب کاربردهای \cmd{sort} معطوف به پردازش داده‌های جدولی است. با استفاده از ترمینولوژی پایگاه‌داده، هر سطر را یک «رکورد» (\en{Record}) و هر بخش از آن سطر را یک «فیلد» (\en{Field}) می‌نامیم (مانند مجوزها، تعداد پیوندها، حجم فایل و نام فایل). ابزار \cmd{sort} توانایی پردازش فیلدهای مجزا را دارد. ما می‌توانیم یک یا چند «فیلدِ کلیدی» (\en{Key Fields}) را به‌عنوان معیار مرتب‌سازی تعیین کنیم. در مثال فوق، با استفاده از گزینه \cmd{-k 5} به برنامه فرمان دادیم که فیلد پنجم (حجم فایل) را به‌عنوان مبنای مرتب‌سازی عددی و معکوس لحاظ کند.

قابلیت \cmd{-k} بسیار منعطف و قدرتمند است، اما پیش از بهره‌گیری از آن، درک منطق تفکیک «فیلدها» (\en{Fields}) در ابزار \cmd{sort} ضروری است. سطر زیر را در نظر بگیرید:

\begin{latin}
\begin{lstlisting}
William Shotts
\end{lstlisting}
\end{latin}

به‌طور پیش‌فرض، \cmd{sort} این سطر را شامل دو فیلد مجزا می‌بیند: فیلد نخست \en{"William"} و فیلد دوم \en{"Shotts"}. در واقع، ابزار \cmd{sort} نویسه‌های \en{Space} و \en{Tab} را به‌عنوان جداکننده‌ی پیش‌فرض در نظر می‌گیرد.

این بدان معناست که کاراکترهای فاصله یا تب به عنوان مرز میان فیلدها عمل کرده و در هنگام عملیات مرتب‌سازی، این جداکننده‌ها در محاسبات مربوط به محتوای هر فیلد لحاظ می‌شوند.

با واکاوی مجدد یک سطر از خروجی دستور \cmd{ls -l} محرز می‌شود که هر رکورد از هشت فیلد مجزا تشکیل شده است که فیلد پنجم آن، اختصاصاً حجم فایل را نمایش می‌دهد:

\begin{latin}
\begin{lstlisting}
-rwxr-xr-x 1 root root 8234216 2016-04-07 17:42 inkscape
\end{lstlisting}
\end{latin}

جهت انجام آزمون‌های عملی پیش‌رو، سناریویی را با استفاده از فایلی حاوی تاریخچه‌ی انتشار سه توزیع پرطرفدار لینوکس در بازه‌ی سال‌های ۲۰۰۶ تا ۲۰۰۸ در نظر می‌گیریم. در این فایل، هر سطر شامل سه فیلدِ نام توزیع (\en{Distribution})، شماره‌ی نسخه (\en{Version}) و تاریخ انتشار (\en{Release Date}) با فرمت \en{MM/DD/YYYY} است:

\begin{latin}
\begin{lstlisting}
SUSE     10.2    12/07/2006
Fedora   10      11/25/2008
SUSE     11.0    06/19/2008
Ubuntu   8.04    04/24/2008
Fedora   8       11/08/2007
SUSE     10.3    10/04/2007
Ubuntu   6.10    10/26/2006
Fedora   7       05/31/2007
Ubuntu   7.10    10/18/2007
Ubuntu   7.04    04/19/2007
SUSE     10.1    05/11/2006
Fedora   6       10/24/2006
Fedora   9       05/13/2008
Ubuntu   6.06    06/01/2006
Ubuntu   8.10    10/30/2008
Fedora   5       03/20/2006
\end{lstlisting}
\end{latin}

با استفاده از یک ویرایشگر متن (مانند \cmd{vim})، داده‌های فوق را وارد کرده و در فایلی با نام \file{distros.txt} ذخیره می‌کنیم. حال با اجرای دستور \cmd{sort} بر روی این فایل، خروجی زیر حاصل می‌شود:

\begin{latin}
\begin{lstlisting}
[me@linuxbox ~]$ sort distros.txt
Fedora   10      11/25/2008
Fedora   5       03/20/2006
Fedora   6       10/24/2006
Fedora   7       05/31/2007
Fedora   8       11/08/2007
Fedora   9       05/13/2008
SUSE     10.1    05/11/2006
SUSE     10.2    12/07/2006
SUSE     10.3    10/04/2007
SUSE     11.0    06/19/2008
Ubuntu   6.06    06/01/2006
Ubuntu   6.10    10/26/2006
Ubuntu   7.04    04/19/2007
Ubuntu   7.10    10/18/2007
Ubuntu   8.04    04/24/2008
Ubuntu   8.10    10/30/2008
\end{lstlisting}
\end{latin}

مشاهده می‌شود که خروجی تا حد زیادی مطلوب است، اما در مرتب‌سازی شماره نسخه‌های توزیع \en{Fedora} اختلالی دیده می‌شود. از آنجا که در مرتب‌سازی پیش‌فرضِ واژه‌نامه‌ای (\en{Lexicographical})، نویسه‌ی «\en{1}» پیش از «\en{5}» قرار می‌گیرد، نسخه‌ی ۱۰ فدورا در ابتدای فهرست و نسخه‌ی ۹ در انتهای آن قرار گرفته است.

برای اصلاح این رفتار، باید از مرتب‌سازی چند‌کلیدی (\en{Multi-key Sorting}) بهره گرفت. هدف ما این است که ابتدا اولویت را به فیلد اول (مرتب‌سازی الفبایی نام توزیع) اختصاص دهیم و سپس فیلد دوم را بر مبنای ارزش عددی مرتب کنیم. ابزار \cmd{sort} امکان فراخوانی مکرر گزینه‌ی \cmd{-k} را برای تعیین چندین «کلیدِ مرتب‌سازی» فراهم می‌کند. شایان ذکر است که هر کلید می‌تواند شامل بازه‌ای از فیلدها باشد؛ اما چنانچه این بازه محدود نشود، \cmd{sort} از فیلد مذکور تا انتهای سطر را به‌عنوان کلید در نظر می‌گیرد. نحوه‌ی پیاده‌سازی این مرتب‌سازی ترکیبی به شرح زیر است:

\begin{latin}
\begin{lstlisting}
[me@linuxbox ~]$ sort --key=1,1 --key=2n distros.txt
Fedora   5       03/20/2006
Fedora   6       10/24/2006
Fedora   7       05/31/2007
Fedora   8       11/08/2007
Fedora   9       05/13/2008
Fedora   10      11/25/2008
SUSE     10.1    05/11/2006
SUSE     10.2    12/07/2006
SUSE     10.3    10/04/2007
SUSE     11.0    06/19/2008
Ubuntu   6.06    06/01/2006
Ubuntu   6.10    10/26/2006
Ubuntu   7.04    04/19/2007
Ubuntu   7.10    10/18/2007
Ubuntu   8.04    04/24/2008
Ubuntu   8.10    10/30/2008
\end{lstlisting}
\end{latin}

اگرچه برای شفافیت بیشتر از فرم طولانی گزینه‌ها استفاده شد، اما عبارت \cmd{-k 1,1 -k 2n} دقیقاً همان عملکرد را دارد. در نخستین فراخوانی کلید، بازه‌ی \cmd{1,1} را تعریف کردیم که به معنای «شروع و پایان در فیلد اول» است؛ این کار باعث می‌شود مرتب‌سازی صرفاً به همان فیلد محدود گردد. در فراخوانی دوم، از \cmd{2n} استفاده کردیم تا فیلد دوم به عنوان مبنای مرتب‌سازی عددی لحاظ شود. شایان ذکر است که می‌توان حروف اختصاری تعیین‌کننده‌ی نوع مرتب‌سازی را مستقیماً به انتهای شماره‌ی فیلد ضمیمه کرد. این حروف همان گزینه‌های عمومی \cmd{sort} هستند؛ مانند \cmd{b} (نادیده گرفتن فواصل ابتدایی)، \cmd{n} (مرتب‌سازی عددی) و \cmd{r} (مرتب‌سازی معکوس).

فیلد سوم در فهرست ما حاوی تاریخی است که ساختار آن برای مرتب‌سازی استاندارد چندان مناسب نیست. در دنیای پردازش داده، معمولاً از فرمت \en{YYYY-MM-DD} استفاده می‌شود تا ترتیب زمانی با ترتیب الفبایی منطبق گردد؛ اما تاریخ‌های فایل ما با فرمت آمریکایی \en{MM/DD/YYYY} درج شده‌اند. برای مرتب‌سازی زمانی چنین فهرستی، ابزار \cmd{sort} قابلیت تعیین «آفست» (\en{Offset}) در داخل فیلد را ارائه می‌دهد:

\begin{latin}
\begin{lstlisting}
[me@linuxbox ~]$ sort -k 3.7nbr -k 3.1nbr -k 3.4nbr distros.txt
Fedora   10      11/25/2008
Ubuntu   8.10    10/30/2008
SUSE     11.0    06/19/2008
Fedora   9       05/13/2008
Ubuntu   8.04    04/24/2008
Fedora   8       11/08/2007
Ubuntu   7.10    10/18/2007
SUSE     10.3    10/04/2007
Fedora   7       05/31/2007
Ubuntu   7.04    04/19/2007
SUSE     10.2    12/07/2006
Ubuntu   6.10    10/26/2006
Fedora   6       10/24/2006
Ubuntu   6.06    06/01/2006
SUSE     10.1    05/11/2006
Fedora   5       03/20/2006
\end{lstlisting}
\end{latin}

در این دستور، با تعیین \cmd{3.7} به برنامه فرمان دادیم که از نویسه‌ی هفتم در فیلد سوم (که آغازگر عدد سال است) برای شروع مرتب‌سازی استفاده کند. به همین ترتیب، آفست‌های \cmd{3.1} و \cmd{3.4} را برای استخراج ماه و روز تعریف کردیم. افزودن گزینه‌های \cmd{n} و \cmd{r} باعث شد تا یک مرتب‌سازی عددی و معکوس (از جدیدترین به قدیمی‌ترین) صورت گیرد. همچنین استفاده از گزینه‌ی \cmd{b} برای حذف فواصل احتمالی ابتدای فیلد الزامی است، چرا که نوسان در تعداد فواصل می‌تواند دقت محاسبات آفست را مختل کند.

در برخی از فایل‌های سیستمی، از نویسه‌هایی به‌جز فاصله (\en{Space}) یا تب (\en{Tab}) برای تفکیک ستون‌ها استفاده می‌شود؛ برای نمونه، ساختار فایل \file{/etc/passwd} را ملاحظه کنید:

\begin{latin}
\begin{lstlisting}
[me@linuxbox ~]$ head /etc/passwd
root:x:0:0:root:/root:/bin/bash
daemon:x:1:1:daemon:/usr/sbin:/bin/sh
bin:x:2:2:bin:/bin:/bin/sh
sys:x:3:3:sys:/dev:/bin/sh
sync:x:4:65534:sync:/bin:/bin/sync
games:x:5:60:games:/usr/games:/bin/sh
man:x:6:12:man:/var/cache/man:/bin/sh
lp:x:7:7:lp:/var/spool/lpd:/bin/sh
mail:x:8:8:mail:/var/mail:/bin/sh
news:x:9:9:news:/var/spool/news:/bin/sh
\end{lstlisting}
\end{latin}

در این فایل، فیلدها توسط نویسه‌ی دونقطه (\cmd{:}) از یکدیگر مجزا شده‌اند. در چنین شرایطی، ابزار \cmd{sort} گزینه‌ی \cmd{-t} (مخفف \en{field-separator}) را جهت تعریف جداکننده‌ی سفارشی ارائه می‌دهد. برای مثال، جهت مرتب‌سازی این فایل بر اساس فیلد هفتم (که نشان‌دهنده‌ی پوسته پیش‌فرض یا \en{Login Shell} کاربر است)، می‌توان از دستور زیر بهره جست:

\begin{latin}
\begin{lstlisting}
[me@linuxbox ~]$ sort -t ':' -k 7 /etc/passwd | head
me:x:1001:1001:Myself,,,:/home/me:/bin/bash
root:x:0:0:root:/root:/bin/bash
dhcp:x:101:102::/nonexistent:/bin/false
gdm:x:106:114:Gnome Display Manager:/var/lib/gdm:/bin/false
hplip:x:104:7:HPLIP system user,,,:/var/run/hplip:/bin/false
klog:x:103:104::/home/klog:/bin/false
messagebus:x:108:119::/var/run/dbus:/bin/false
polkituser:x:110:122:PolicyKit,,,:/var/run/PolicyKit:/bin/false
pulse:x:107:116:PulseAudio daemon,,,:/var/run/pulse:/bin/false
\end{lstlisting}
\end{latin}

با تعیین صریح نویسه‌ی دونقطه به‌عنوان جداکننده‌ی فیلد (\en{Field Delimiter})، امکان مرتب‌سازی دقیق داده‌ها بر مبنای ستون هفتم فراهم شده است.

\section{دستور uniq}

در مقایسه با ابزار \cmd{sort} دستور \cmd{uniq} ساختار ساده‌تری دارد و وظیفه‌ای مشخص را دنبال می‌کند: این برنامه با دریافت یک فایل مرتب‌سازی‌شده (یا ورودی استاندارد)، سطرهای تکراری متوالی را حذف کرده و نتیجه را در خروجی استاندارد چاپ می‌کند. از این رو، \cmd{uniq} غالباً در ترکیب با \cmd{sort} جهت پاک‌سازی خروجی از رکوردهای تکراری به‌کار می‌رود.

\begin{note}
\textbf{نکته:} اگرچه \cmd{uniq} ابزاری کلاسیک در یونیکس محسوب می‌شود، اما نسخه‌ی \en{GNU} ابزار \cmd{sort} دارای گزینه \cmd{-u} است که می‌تواند هم‌زمان با مرتب‌سازی، سطرهای تکراری را نیز فیلتر کند.
\end{note}

برای درک بهتر عملکرد این ابزار، فایلی آزمایشی ایجاد می‌کنیم:

\begin{latin}
\begin{lstlisting}
[me@linuxbox ~]$ cat > foo.txt
a
b
c
a
b
c
\end{lstlisting}
\end{latin}

پس از تایپ محتوا، برای خروج از وضعیت ورود داده، کلید \cmd{Ctrl+d} را فشار دهید. حال اگر \cmd{uniq} را مستقیماً روی این فایل اجرا کنیم، نتیجه به شکل زیر خواهد بود:

\begin{latin}
\begin{lstlisting}
[me@linuxbox ~]$ uniq foo.txt
a
b
c
a
b
c
\end{lstlisting}
\end{latin}

مشاهده می‌کنید که تغییری در خروجی حاصل نشده و سطرهای تکراری حذف نشده‌اند. برای اینکه \cmd{uniq} بتواند وظیفه‌ی خود را به‌درستی انجام دهد، ورودی آن لزوماً باید از قبل مرتب‌سازی شده باشد:

\begin{latin}
\begin{lstlisting}
[me@linuxbox ~]$ sort foo.txt | uniq
a
b
c
\end{lstlisting}
\end{latin}

دلیل این رفتار آن است که \cmd{uniq} صرفاً سطرهای تکراری که بلافاصله پس از یکدیگر (\en{Adjacent}) قرار دارند را شناسایی و حذف می‌کند.

ابزار \cmd{uniq} دارای چندین گزینه کاربردی است که امکان مدیریت دقیق‌تر داده‌های تکراری را فراهم می‌کند. جدول زیر، متداول‌ترین گزینه‌های این دستور را معرفی می‌کند:

\begin{table}[htbp]
\centering
\caption{گزینه‌های کاربردی دستور \cmd{uniq}}
\label{tab:uniq-options}
\begin{tabularx}{\textwidth}{l l X}
\toprule
\textbf{گزینه کوتاه} & \textbf{گزینه بلند} & \textbf{شرح عملکرد} \\
\midrule
\cmd{-c} & \cmd{--count} & شمارش تعداد تکرار هر خط را محاسبه کرده و آن را به‌عنوان یک پیشوند عددی در ابتدای خط خروجی درج می‌کند. \\
\addlinespace
\cmd{-d} & \cmd{--repeated} & تنها خطوطی را در خروجی نمایش می‌دهد که حداقل یک بار تکرار شده‌اند. به‌ازای هر گروه از خطوط تکراری، فقط یک نمونه چاپ می‌شود. \\
\addlinespace
\cmd{-f n} & \cmd{--skip-fields=n} & پیش از شروع مقایسه، تعداد \en{n} فیلد ابتدایی هر سطر را نادیده می‌گیرد. فیلدها دنباله‌ای از کاراکترهای غیرفاصله هستند که با فاصله یا تب از هم جدا می‌شوند. \\
\addlinespace
\cmd{-i} & \cmd{--ignore-case} & در فرآیند مقایسه خطوط، تفاوت میان حروف بزرگ و کوچک را نادیده می‌گیرد. \\
\addlinespace
\cmd{-s n} & \cmd{--skip-chars=n} & پیش از شروع مقایسه، تعداد \en{n} کاراکتر نخست هر سطر را نادیده می‌گیرد. \\
\addlinespace
\cmd{-u} & \cmd{--unique} & تنها خطوطی را در خروجی نمایش می‌دهد که در کل ورودی، تکرار نشده و منحصربه‌فرد هستند. \\
\bottomrule
\end{tabularx}
\end{table}

در این مثال، ابزار \cmd{uniq} با استفاده از گزینه‌ی \cmd{-c} برای گزارش تعداد فراوانی هر سطر در فایل متنی به‌کار گرفته شده است:

\begin{latin}
\begin{lstlisting}
[me@linuxbox ~]$ sort foo.txt | uniq -c
     2 a
     2 b
     2 c
\end{lstlisting}
\end{latin}

در خروجی فوق، عددی که در ابتدای هر سطر درج شده، نشان‌دهنده‌ی تعداد دفعات تکرار آن رکورد در فایل اصلی است. این ترکیب (\cmd{sort | uniq -c}) یکی از پرکاربردترین الگوها در خط فرمان لینوکس برای استخراج آمار و تحلیل توزیع داده‌ها محسوب می‌شود.

\section{ابزارهای برش و تفکیک متن}

سه برنامه‌ای که در ادامه بررسی خواهیم کرد، جهت استخراج ستون‌های متنی از فایل‌ها و بازترکیب آن‌ها به شیوه‌های کاربردی طراحی شده‌اند.

\subsection{دستور cut}

برنامه‌ی \cmd{cut} برای استخراج بخش‌های مشخصی از یک سطر متن و ارسال آن به خروجی استاندارد (\en{stdout}) به‌کار می‌رود. این ابزار قابلیت پذیرش چندین فایل ورودی را دارد و همچنین می‌تواند داده‌ها را از ورودی استاندارد (\en{stdin}) دریافت کند.

تعیین محدوده‌ی استخراج در هر سطر با استفاده از گزینه‌های جدول زیر انجام می‌شود که پارامترهای نسبتاً متنوعی را در اختیار کاربر قرار می‌دهد:

\begin{table}[htbp]
\centering
\caption{گزینه‌های کاربردی دستور \cmd{cut}}
\label{tab:cut-options}
\begin{tabularx}{\textwidth}{l l X}
\toprule
\textbf{گزینه کوتاه} & \textbf{گزینه بلند} & \textbf{شرح عملکرد} \\
\midrule
\cmd{-c list} & \cmd{--characters=list} & استخراج کاراکترهای مشخص‌شده در \file{list} از هر سطر. \file{list} شامل یک یا چند عدد یا بازه (مانند \cmd{1-5} یا \cmd{3-7,9}) است که با کاما از هم جدا می‌شوند. شماره‌گذاری کاراکترها از ۱ آغاز می‌شود. \\
\addlinespace
\cmd{-f list} & \cmd{--fields=list} & استخراج فیلدهای مشخص‌شده در \file{list} از هر سطر. فیلدها به‌طور پیش‌فرض با نویسهٔ \en{Tab} از هم جدا می‌شوند. هر سطری که فاقد جداکننده باشد، به‌طور کامل چاپ می‌شود مگر اینکه گزینهٔ \cmd{-s} نیز به کار رود. \\
\addlinespace
\cmd{-d delim} & \cmd{--delimiter=delim} & تعیین نویسهٔ \file{delim} به‌عنوان جداکنندهٔ فیلدها هنگام استفاده از گزینهٔ \cmd{-f}. در صورت عدم استفاده از این گزینه، نویسهٔ \en{Tab} به‌عنوان جداکنندهٔ پیش‌فرض در نظر گرفته می‌شود. \\
\addlinespace
\cmd{--complement} & --- & معکوس‌سازی عملیات انتخاب؛ به‌جای استخراج بخش‌های مشخص‌شده توسط \cmd{-c} یا \cmd{-f}، تمامی بخش‌های دیگر سطر استخراج می‌شوند. این گزینه یک افزونهٔ \en{GNU} است. \\
\bottomrule
\end{tabularx}
\end{table}

همان‌طور که ملاحظه می‌شود، مکانیزم استخراج متن در ابزار \cmd{cut} تا حدودی محدود است. این برنامه بیشترین کارایی را در پردازش خروجی‌های تولیدشده توسط سایر برنامه‌ها (داده‌های ساختاریافته) دارد و برای متونی که مستقیماً توسط انسان تایپ شده‌اند (و ممکن است نظم فیلدبندی دقیقی نداشته باشند)، گزینه‌ی ایده‌آلی محسوب نمی‌شود.

به منظور اطمینان از ساختار یافته بودن فایل \file{distros.txt} برای استفاده در مثال‌های ابزار \cmd{cut} سطرها را واکاوی می‌کنیم. با به‌کارگیری دستور \cmd{cat} همراه با گزینه \cmd{-A} می‌توان وجود نویسه‌های کنترلی نظیر \en{Tab} را احراز کرد:

\begin{latin}
\begin{lstlisting}
[me@linuxbox ~]$ cat -A distros.txt
SUSE^I10.2^I12/07/2006$
Fedora^I10^I11/25/2008$
SUSE^I11.0^I06/19/2008$
Ubuntu^I8.04^I04/24/2008$
Fedora^I8^I11/08/2007$
SUSE^I10.3^I10/04/2007$
Ubuntu^I6.10^I10/26/2006$
Fedora^I7^I05/31/2007$
Ubuntu^I7.10^I10/18/2007$
Ubuntu^I7.04^I04/19/2007$
SUSE^I10.1^I05/11/2006$
Fedora^I6^I10/24/2006$
Fedora^I9^I05/13/2008$
Ubuntu^I6.06^I06/01/2006$
Ubuntu^I8.10^I10/30/2008$
Fedora^I5^I03/20/2006$
\end{lstlisting}
\end{latin}

همان‌طور که مشاهده می‌شود، داده‌ها کاملاً منسجم هستند؛ فواصل اضافی در متن وجود ندارد و فیلدها صرفاً با یک نویسه‌ی \en{Tab} (که در اینجا با \cmd{\^{}I} نمایش داده شده) از هم تفکیک شده‌اند. با توجه به اینکه \en{Tab} جداکننده‌ی پیش‌فرض در ابزار \cmd{cut} است، استفاده از گزینه‌ی \cmd{-f} (مخفف \en{fields}) بهینه‌ترین روش برای استخراج ستون‌ها خواهد بود:

\begin{latin}
\begin{lstlisting}
[me@linuxbox ~]$ cut -f 3 distros.txt
12/07/2006
11/25/2008
06/19/2008
04/24/2008
11/08/2007
10/04/2007
10/26/2006
05/31/2007
10/18/2007
04/19/2007
05/11/2006
10/24/2006
05/13/2008
06/01/2006
10/30/2008
03/20/2006
\end{lstlisting}
\end{latin}

در فایل‌هایی که با \en{Tab} پیکربندی شده‌اند، استخراج بر اساس فیلد بسیار مطمئن‌تر از استخراج بر اساس کاراکتر است؛ چرا که طول متغیر رشته‌ها در هر سطر، محاسبه‌ی موقعیت دقیق کاراکترها را دشوار می‌سازد. با این حال، در مثال فوق، از آنجا که داده‌های فیلد سوم (تاریخ) دارای طول ثابتی هستند، می‌توانیم از استخراج کاراکتری برای جداسازی بخش «سال» استفاده کنیم:

\begin{latin}
\begin{lstlisting}
[me@linuxbox ~]$ cut -f 3 distros.txt | cut -c 7-10
2006
2008
2008
2008
2007
2007
2006
2007
2007
2007
2006
2006
2008
2006
2008
2006
\end{lstlisting}
\end{latin}

در این فرآیند، خروجیِ دستور اول به یک دستور \cmd{cut} دیگر پایپ (\en{Pipe}) شده است تا صرفاً کاراکترهای موقعیت ۷ تا ۱۰ (محدوده‌ی سال) استخراج شوند. عبارت \cmd{7-10} نمونه‌ای از تعیین «بازه» (\en{Range}) در خط فرمان است. برای مطالعه‌ی شیوه‌های پیچیده‌ترِ تعیین بازه، می‌توانید به مستندات سیستم (\cmd{man cut}) مراجعه کنید.

\subsection{تبدیل نویسه‌های Tab به Space}

فایل \file{distros.txt} به دلیل استفاده از نویسه‌ی \en{Tab}، ساختاری ایده‌آل برای استخراج فیلدها با ابزار \cmd{cut} دارد. اما چنانچه بخواهیم فایلی را بر اساس موقعیت دقیق کاراکترها (و نه لزوماً فیلدها) پردازش کنیم، وجود تب‌ها چالش‌برانگیز خواهد بود؛ زیرا هر کاراکتر تب در فایل تنها یک بایت اشغال می‌کند اما در نمایشگر معادل چندین فاصله دیده می‌شود. برای رفع این مشکل، باید تب‌ها را با تعداد معادلِ فاصله (\en{Space}) جایگزین کنیم.

بسته‌ی \en{GNU Coreutils} برای این منظور ابزار \cmd{expand} را ارائه می‌دهد. این برنامه با دریافت فایل ورودی یا خواندن از ورودی استاندارد، تب‌ها را به فاصله تبدیل کرده و نتیجه را در خروجی استاندارد چاپ می‌کند.

با پردازش فایل \file{distros.txt} توسط دستور \cmd{expand} می‌توانیم به‌راحتی از \cmd{cut -c} برای استخراج هر بازه‌ی دلخواهی از کاراکترها استفاده کنیم. برای نمونه، جهت استخراج سال انتشار، ابتدا فایل را \cmd{expand} کرده و سپس تمامی کاراکترهای موجود از موقعیت ۲۳ تا انتهای سطر را جدا می‌کنیم:

\begin{latin}
\begin{lstlisting}
[me@linuxbox ~]$ expand distros.txt | cut -c 23-
\end{lstlisting}
\end{latin}

همچنین ابزار معکوس این برنامه، یعنی \cmd{unexpand} نیز در دسترس است که فاصله‌ها را (به‌طور پیش‌فرض در ابتدای خطوط) به نویسه‌ی \en{Tab} تبدیل می‌کند. این ابزار بیشتر برای کاهش حجم فایل‌ها از طریق جایگزینی فاصله‌ها با نویسه‌ی فشرده‌تر \en{Tab} کاربرد دارد. به‌طور پیش‌فرض، \cmd{unexpand} تنها فاصله‌های ابتدای هر خط را تبدیل می‌کند، اما با گزینه‌ی \cmd{-a} می‌توان تمام فاصله‌ها را به تب تبدیل نمود.

در هنگام کار با فیلدها، الزامی به استفاده از نویسه‌ی \en{Tab} وجود ندارد. ابزار \cmd{cut} با استفاده از گزینه‌ی \cmd{-d} امکان تعیین هر کاراکتر دلخواهی را به‌عنوان جداکننده فراهم می‌کند. در مثال زیر، فیلد اول از فایل \file{/etc/passwd} (که حاوی نام‌های کاربری است) را استخراج می‌کنیم:

\begin{latin}
\begin{lstlisting}
[me@linuxbox ~]$ cut -d ':' -f 1 /etc/passwd | head
root
daemon
bin
sys
sync
games
man
lp
mail
news
\end{lstlisting}
\end{latin}

با استفاده از پارامتر \cmd{-d ':'} به برنامه اعلام کردیم که در این فایل، نویسه‌ی دونقطه نقش جداکننده‌ی ستون‌ها را بر عهده دارد.

\subsection{دستور paste}

دستور \cmd{paste} عملکردی معکوس نسبت به ابزار \cmd{cut} دارد. این برنامه به‌جای جداسازی یک ستون متنی از فایل، یک یا چند ستون را به آن الحاق می‌کند. مکانیزم کارکرد آن بدین‌صورت است که چندین فایل را به‌صورت موازی خوانده و فیلدهای متناظر در هر یک را با هم ترکیب کرده و در قالب یک جریان واحد به خروجی استاندارد ارسال می‌کند.

مشابه \cmd{cut} ابزار \cmd{paste} نیز آرگومان‌های متعددی را در قالب فایل یا ورودی استاندارد پذیراست. جهت نمایش کاربرد عملی این ابزار، فایل \file{distros.txt} را به‌گونه‌ای بازآرایی می‌کنیم که فهرستی مبتنی بر توالیِ زمانی از نسخه‌ها حاصل شود.

ابتدا با بهره‌گیری از آموخته‌های پیشین درباره دستور \cmd{sort} فهرستی از توزیع‌ها را بر اساس تاریخ مرتب کرده و خروجی را در فایلی به نام \file{distros-by-date.txt} ذخیره می‌کنیم:

\begin{latin}
\begin{lstlisting}
[me@linuxbox ~]$ sort -k 3.7nbr -k 3.1nbr -k 3.4nbr distros.txt > distros-by-date.txt
\end{lstlisting}
\end{latin}

در گام بعد، با استفاده از دستور \cmd{cut} دو فیلد نخست (نام توزیع و شماره نسخه) را استخراج کرده و در فایل \file{distros-versions.txt} قرار می‌دهیم:

\begin{latin}
\begin{lstlisting}
[me@linuxbox ~]$ cut -f 1,2 distros-by-date.txt > distros-versions.txt
[me@linuxbox ~]$ head distros-versions.txt
Fedora   10
Ubuntu   8.10
SUSE     11.0
Fedora   9
Ubuntu   8.04
Fedora   8
Ubuntu   7.10
SUSE     10.3
Fedora   7
Ubuntu   7.04
\end{lstlisting}
\end{latin}

به‌عنوان آخرین مرحله از فرآیند آماده‌سازی، تاریخ‌های انتشار را نیز استخراج کرده و در فایلی با نام \file{distros-dates.txt} ذخیره می‌کنیم:

\begin{latin}
\begin{lstlisting}
[me@linuxbox ~]$ cut -f 3 distros-by-date.txt > distros-dates.txt
[me@linuxbox ~]$ head distros-dates.txt
11/25/2008
10/30/2008
06/19/2008
05/13/2008
04/24/2008
11/08/2007
10/18/2007
10/04/2007
05/31/2007
04/19/2007
\end{lstlisting}
\end{latin}

اکنون که تمامی اجزای مورد نیاز تفکیک شده‌اند، با استفاده از دستور \cmd{paste} ستون تاریخ‌ها را پیش از نام و نسخه‌ی توزیع‌ها قرار می‌دهیم تا یک فهرست ترتیبی (بر اساس زمان) ایجاد شود. این عملیات با تنظیم ترتیب آرگومان‌های ورودی در دستور \cmd{paste} به‌سادگی امکان‌پذیر است:

\begin{latin}
\begin{lstlisting}
[me@linuxbox ~]$ paste distros-dates.txt distros-versions.txt
11/25/2008	Fedora	10
10/30/2008	Ubuntu	8.10
06/19/2008	SUSE	      11.0
05/13/2008	Fedora	9
04/24/2008	Ubuntu	8.04
11/08/2007	Fedora	8
10/18/2007	Ubuntu	7.10
10/04/2007	SUSE	      10.3
05/31/2007	Fedora	7
04/19/2007	Ubuntu	7.04
12/07/2006	SUSE	      10.2
10/26/2006	Ubuntu	6.10
10/24/2006	Fedora	6
06/01/2006	Ubuntu	6.06
05/11/2006	SUSE	      10.1
03/20/2006	Fedora      5
\end{lstlisting}
\end{latin}

\subsection{دستور join}

از برخی جنبه‌ها، \cmd{join} شباهت‌هایی به ابزار \cmd{paste} دارد، چرا که ستون‌های جدیدی را به فایل اضافه می‌کند؛ اما مکانیزم عملکرد آن کاملاً منحصربه‌فرد است. عملیات \cmd{join} به‌طور معمول در ادبیات سیستم‌های مدیریت پایگاه داده رابطه‌ای (\en{Relational Database Management System - RDBMS}) ریشه دارد؛ جایی که داده‌ها از چندین جدول مختلف، بر اساس یک فیلد کلید مشترک (\en{Common Key Field}) با یکدیگر ترکیب می‌شوند تا خروجی مورد نظر حاصل گردد.

برنامه‌ی \cmd{join} دقیقاً همین وظیفه را در سطح فایل‌های متنی بر عهده دارد. این برنامه داده‌های موجود در دو فایل را بر اساس یک فیلد کلید یکسان به یکدیگر پیوند (\en{Link}) می‌دهد.

برای درک بهتر کاربرد عملیات \cmd{join} در یک پایگاه داده رابطه‌ای، فرض کنید یک پایگاه داده کوچک متشکل از دو جدول داریم که هر کدام حاوی یک رکورد هستند. جدول نخست تحت عنوان \en{CUSTOMERS} دارای سه فیلد است: شماره مشتری (\en{CUSTNUM})، نام (\en{FNAME}) و نام خانوادگی (\en{LNAME}):

\begin{latin}
\begin{lstlisting}
CUSTNUM	FNAME	LNAME
4681934	John	Smith
\end{lstlisting}
\end{latin}

جدول دوم با نام \en{ORDERS} شامل چهار فیلد است: شماره سفارش (\en{ORDERNUM})، شماره مشتری (\en{CUSTNUM})، تعداد (\en{QUAN}) و شرح کالا (\en{ITEM}):

\begin{latin}
\begin{lstlisting}
ITEM	QUAN	CUSTNUM	ORDERNUM
Blue Widget	1	4681934	3014953305
\end{lstlisting}
\end{latin}

شایان ذکر است که هر دو جدول دارای فیلد مشترک \en{CUSTNUM} هستند. این اشتراک نقشی کلیدی ایفا می‌کند، زیرا امکان برقراری رابطه (\en{Relationship}) میان دو جدول را میسر می‌سازد.

اجرای عملیات \cmd{join} به ما اجازه می‌دهد فیلدهای این دو جدول را برای دستیابی به خروجی کاربردی‌تر (مثلاً صدور یک فاکتور خرید) با هم ادغام کنیم. با انطباق مقادیر موجود در فیلد \en{CUSTNUM} در هر دو جدول، دستور \cmd{join} می‌تواند نتیجه‌ی ترکیبی زیر را ارائه دهد:

\begin{latin}
\begin{lstlisting}
ITEM	QUAN	FNAME	LNAME
Blue Widget	1	John	Smith
\end{lstlisting}
\end{latin}

به منظور تشریح عملکرد ابزار \cmd{join} نیازمند ایجاد فایل‌هایی با یک «فیلد کلید» (\en{Common Key}) مشترک هستیم. در این سناریو، از فایل \file{distros-by-date.txt} به عنوان منبع اصلی استفاده کرده و دو فایل مجزا از آن مشتق می‌کنیم. فایل نخست شامل تاریخ انتشار (که نقش کلید مشترک را ایفا می‌کند) و نام توزیع خواهد بود:

\begin{latin}
\begin{lstlisting}
[me@linuxbox ~]$ cut -f 1,1 distros-by-date.txt > distros-names.txt
[me@linuxbox ~]$ paste distros-dates.txt distros-names.txt > distros-key-names.txt
[me@linuxbox ~]$ head distros-key-names.txt
11/25/2008	Fedora
10/30/2008	Ubuntu
06/19/2008	SUSE
05/13/2008	Fedora
04/24/2008	Ubuntu
11/08/2007	Fedora
10/18/2007	Ubuntu
10/04/2007	SUSE
05/31/2007	Fedora
04/19/2007	Ubuntu
\end{lstlisting}
\end{latin}

فایل دوم را به گونه‌ای تنظیم می‌کنیم که شامل تاریخ انتشار و شماره نسخه توزیع‌ها باشد:

\begin{latin}
\begin{lstlisting}
[me@linuxbox ~]$ cut -f 2,2 distros-by-date.txt > distros-vernums.txt
[me@linuxbox ~]$ paste distros-dates.txt distros-vernums.txt > distros-key-vernums.txt
[me@linuxbox ~]$ head distros-key-vernums.txt
11/25/2008	10
10/30/2008	8.10
06/19/2008	11.0
05/13/2008	9
04/24/2008	8.04
11/08/2007	8
10/18/2007	7.10
10/04/2007	10.3
05/31/2007	7
04/19/2007	7.04
\end{lstlisting}
\end{latin}

اکنون دو منبع داده در اختیار داریم که توسط فیلد «تاریخ انتشار» به یکدیگر مرتبط شده‌اند. ذکر این نکته ضروری است که برای عملکرد صحیح ابزار \cmd{join} فایل‌های ورودی الزاماً باید بر اساس فیلد کلید مرتب‌سازی (\en{Sort}) شده باشند.

\begin{latin}
\begin{lstlisting}
[me@linuxbox ~]$ join distros-key-names.txt distros-key-vernums.txt | head
11/25/2008 Fedora 10
10/30/2008 Ubuntu 8.10
06/19/2008 SUSE 11.0
05/13/2008 Fedora 9
04/24/2008 Ubuntu 8.04
11/08/2007 Fedora 8
10/18/2007 Ubuntu 7.10
10/04/2007 SUSE 10.3
05/31/2007 Fedora 7
04/19/2007 Ubuntu 7.04
\end{lstlisting}
\end{latin}

شایان ذکر است که ابزار \cmd{join} به‌صورت پیش‌فرض از نویسه‌ی فاصله (\en{White Space}) به عنوان جداکننده‌ی فیلدهای ورودی و خروجی استفاده می‌کند. این رفتار استاندارد از طریق گزینه‌های برنامه قابل سفارشی‌سازی است. برای بررسی پارامترهای پیشرفته‌تر، می‌توانید به مستندات سیستم (\cmd{man join}) مراجعه نمایید.

\subsection{دستور tac}

دستور \cmd{tac} از نظر عملکردی مشابه \cmd{cat} است، با این تفاوت که پردازش سطرها را به‌صورت معکوس (از انتها به ابتدا) انجام می‌دهد. نام این ابزار نیز دقیقاً برعکس‌شده‌ی کلمه‌ی \en{cat} است تا ماهیت معکوس آن را تداعی کند. یکی از کاربردهای رایج \cmd{tac} در مدیریت سیستم، بازآرایی فایل‌های گزارش (\en{Log Files}) برای بررسی سریع‌تر وقایع اخیر است.

به‌طور معمول، سیستم‌های لینوکسی وقایع جدید را به انتهای فایل‌های گزارش ضمیمه می‌کنند، به گونه‌ای که قدیمی‌ترین رکورد در ابتدای فایل و جدیدترین آن‌ها در انتها قرار می‌گیرد:

\begin{latin}
\begin{lstlisting}
[me@linuxbox ~]$ tail /var/log/backup.log
Wed Aug 14 06:00:32 AM EDT 2025: backup (0.6) of linuxbox started.
Wed Aug 14 06:08:01 AM EDT 2025: backup of linuxbox finished.

Thu Aug 15 07:36:30 AM EDT 2025: backup (0.6) of linuxbox started.
Thu Aug 15 07:45:57 AM EDT 2025: backup of linuxbox finished.

Fri Aug 16 07:37:57 AM EDT 2025: backup (0.6) of linuxbox started.
Fri Aug 16 07:44:11 AM EDT 2025: backup of linuxbox finished.
\end{lstlisting}
\end{latin}

چنانچه تحلیل‌گر نیاز داشته باشد تازه‌ترین رویدادها را در صدر فهرست مشاهده کند، می‌تواند خروجی را به ابزار \cmd{tac} پایپ (\en{Pipe}) کند:

\begin{latin}
\begin{lstlisting}
[me@linuxbox ~]$ tail /var/log/backup.log | tac
Fri Aug 16 07:44:11 AM EDT 2025: backup of linuxbox finished.
Fri Aug 16 07:37:57 AM EDT 2025: backup (0.6) of linuxbox started.

Thu Aug 15 07:45:57 AM EDT 2025: backup of linuxbox finished.
Thu Aug 15 07:36:30 AM EDT 2025: backup (0.6) of linuxbox started.

Wed Aug 14 06:08:01 AM EDT 2025: backup of linuxbox finished.
Wed Aug 14 06:00:32 AM EDT 2025: backup (0.6) of linuxbox started.
\end{lstlisting}
\end{latin}

\subsection{دستور rev}

دستور \cmd{rev} عملکردی مشابه \cmd{tac} دارد، اما به‌جای معکوس کردن ترتیب سطرها، ترتیب نویسه‌ها را در هر سطر وارونه می‌کند:

\begin{latin}
\begin{lstlisting}
[me@linuxbox ~]$ echo "This is a test." | rev
.tset a si sihT
\end{lstlisting}
\end{latin}

اگرچه این قابلیت در ظاهر صرفاً جنبه‌ی نمایشی دارد، اما در ترکیب با ابزارهایی نظیر \cmd{cut} قدرت واقعی خود را در پردازش متون با طول متغیر نشان می‌دهد. بار دیگر فایل \file{backup.log} را در نظر بگیرید؛ فرض کنید قصد داریم علامت نقطه (\cmd{.}) را از انتهای تمامی سطرها حذف کنیم. از آنجا که طول سطرها یکسان نیست، نمی‌توان با استفاده از ابزار \cmd{cut} موقعیت ثابتی را از انتهای سطر هدف قرار داد. اینجاست که \cmd{rev} راهکاری هوشمندانه ارائه می‌دهد: با معکوس کردن سطر، نویسه‌ی انتهایی به ابتدای سطر منتقل می‌شود و به‌راحتی توسط \cmd{cut} قابل حذف خواهد بود.

\begin{latin}
\begin{lstlisting}
[me@linuxbox ~]$ tail /var/log/backup.log | rev | cut -c 2- | rev
Wed Aug 14 06:00:32 AM EDT 2025: backup (0.6) of linuxbox started
Wed Aug 14 06:08:01 AM EDT 2025: backup of linuxbox finished

Thu Aug 15 07:36:30 AM EDT 2025: backup (0.6) of linuxbox started
Thu Aug 15 07:45:57 AM EDT 2025: backup of linuxbox finished

Fri Aug 16 07:37:57 AM EDT 2025: backup (0.6) of linuxbox started
Fri Aug 16 07:44:11 AM EDT 2025: backup of linuxbox finished
\end{lstlisting}
\end{latin}

در این زنجیره‌ی دستوری (\en{Pipeline})، ابتدا توسط \cmd{rev} محتوای هر سطر واژگون می‌شود. سپس با \cmd{cut -c 2-} تمامی نویسه‌ها از موقعیت دوم به بعد (یعنی حذف اولین کاراکتر که همان نقطه‌ی انتهایی سابق است) استخراج می‌گردند. در نهایت، با فراخوانی مجدد \cmd{rev} متن به حالت طبیعی و خوانا بازگردانده می‌شود.

\section{ابزارهای مقایسه متن}

در بسیاری از سناریوها، مقایسه‌ی نسخه‌های مختلف فایل‌های متنی ضرورتی اجتناب‌ناپذیر است. این موضوع برای مدیران سیستم و توسعه‌دهندگان نرم‌افزار اهمیت ویژه‌ای دارد؛ برای نمونه، یک مدیر سیستم ممکن است برای ریشه‌یابی اختلالات، فایل پیکربندی (\en{Config}) فعلی را با نسخه‌ی پشتیبان قبلی مقایسه کند، یا یک برنامه‌نویس بخواهد تغییرات اعمال‌شده در کد منبع را در طول زمان رصد نماید.

\subsection{دستور comm}

برنامه‌ی \cmd{comm} دو فایل متنی را سطر به سطر مقایسه کرده و سطرهای منحصربه‌فرد هر فایل و همچنین سطرهای مشترک را تفکیک می‌کند. برای درک بهتر، دو فایل تقریباً مشابه ایجاد می‌کنیم:

\begin{latin}
\begin{lstlisting}
[me@linuxbox ~]$ cat > file1.txt
a
b
c
d
[me@linuxbox ~]$ cat > file2.txt
b
c
d
e
\end{lstlisting}
\end{latin}

سپس با استفاده از دستور \cmd{comm} آن‌ها را با هم مقایسه می‌کنیم:

\begin{latin}
\begin{lstlisting}
[me@linuxbox ~]$ comm file1.txt file2.txt
a
       b
       c
       d
   e
\end{lstlisting}
\end{latin}

همان‌طور که ملاحظه می‌شود، خروجی \cmd{comm} در سه ستون دسته‌بندی می‌شود:

\begin{itemize}
  \item \textbf{ستون اول:} سطرهایی که صرفاً در فایل اول موجودند.
  \item \textbf{ستون دوم:} سطرهایی که صرفاً در فایل دوم یافت می‌شوند.
  \item \textbf{ستون سوم:} سطرهایی که در هر دو فایل مشترک هستند.
\end{itemize}

ابزار \cmd{comm} از گزینه‌هایی به فرم \cmd{-n} بهره می‌برد که در آن \en{n} می‌تواند \cmd{1}، \cmd{2} یا \cmd{3} باشد. این گزینه‌ها برای مخفی‌سازی ستون‌های متناظر به‌کار می‌روند. برای مثال، اگر بخواهیم صرفاً سطرهای مشترک میان دو فایل را استخراج کنیم، باید ستون‌های اول و دوم را حذف نماییم:

\begin{latin}
\begin{lstlisting}
[me@linuxbox ~]$ comm -12 file1.txt file2.txt
b
c
d
\end{lstlisting}
\end{latin}

توجه داشته باشید که ابزار \cmd{comm} برای عملکرد صحیح، نیازمند آن است که هر دو فایل ورودی از پیش مرتب‌سازی شده باشند.

\subsection{دستور diff}

مشابه برنامه‌ی \cmd{comm} ابزار \cmd{diff} نیز برای شناسایی تفاوت‌های میان فایل‌ها به کار می‌رود. با این حال، \cmd{diff} ابزاری به‌مراتب توانمندتر و پیچیده‌تر است که از قالب‌های خروجی متنوعی پشتیبانی می‌کند و قابلیت پردازش هم‌زمان مجموعه‌های حجیمی از فایل‌های متنی را دارد. توسعه‌دهندگان نرم‌افزار به‌وفور از \cmd{diff} برای واکاوی تغییرات در نسخه‌های مختلف کد منبع استفاده می‌کنند؛ از این رو، این ابزار قابلیت پیمایش بازگشتی (\en{recursive}) دایرکتوری‌های پروژه را که در اصطلاح تخصصی «ساختار درختی کد منبع» (\en{Source Tree}) نامیده می‌شوند، داراست.

یکی از کاربردهای کلیدی \cmd{diff} ایجاد فایل‌های وصله (\en{Patch}) است. این فایل‌ها توسط برنامه‌هایی نظیر \cmd{patch} (که در ادامه به آن خواهیم پرداخت) به‌منظور تبدیل یک نسخه از فایل به نسخه‌ای دیگر مورد استفاده قرار می‌گیرند.

چنانچه از \cmd{diff} برای مقایسه‌ی فایل‌های نمونه‌ی قبلی استفاده کنیم، خروجی به شرح زیر خواهد بود:

\begin{latin}
\begin{lstlisting}
[me@linuxbox ~]$ diff file1.txt file2.txt
1d0
< a
4a4
> e
\end{lstlisting}
\end{latin}

در اینجا فرمت پیش‌فرض خروجی را مشاهده می‌کنید که گزارشی فشرده از تفاوت‌های میان دو فایل ارائه می‌دهد. در این قالب، هر گروه از تغییرات با یک «دستور ویرایشی» (\en{Change Command}) همراه است. این دستورات به‌صورت «محدوده، نوع عملیات، محدوده» بیان می‌شوند تا موقعیت و ماهیت تغییرات لازم برای تبدیل فایل اول به فایل دوم را مطابق جدول زیر تبیین کنند.

\begin{table}[htbp]
\centering
\caption{دستورات ویرایشی در خروجی \cmd{diff}}
\label{tab:diff-commands}
\begin{tabularx}{\textwidth}{l X}
\toprule
\textbf{دستور} & \textbf{شرح عملکرد} \\
\midrule
\cmd{r1a r2} & نشان‌دهنده‌ی عملیات افزودن (\en{Add}) است. این دستور به این معناست که برای تبدیل فایل اول به فایل دوم، باید سطرهای واقع در محدوده‌ی \en{r2} از فایل دوم را، بلافاصله پس از سطر \en{r1} در فایل اول، اضافه کرد. در اینجا، \en{r1} به سطری در فایل اول اشاره دارد که محل درج تغییرات جدید پس از آن است، و \en{r2} محدوده‌ی سطرهایی از فایل دوم است که باید افزوده شوند. \\
\addlinespace
\cmd{r1c r2} & نشان‌دهنده‌ی عملیات تغییر (\en{Change}) است. این دستور بیان می‌کند که سطرهای واقع در محدوده‌ی \en{r1} از فایل اول باید با سطرهای واقع در محدوده‌ی \en{r2} از فایل دوم جایگزین شوند. از نظر عملکردی، این دستور معادل انجام همزمان یک عملیات حذف و سپس افزودن است، با این تفاوت که \cmd{diff} آن را در قالبی فشرده‌تر و خواناتر گزارش می‌دهد. \\
\addlinespace
\cmd{r1d r2} & نشان‌دهنده‌ی عملیات حذف (\en{Delete}) است. این دستور به این معناست که سطرهای واقع در محدوده‌ی \en{r1} باید از فایل اول حذف شوند تا فایل اول به فایل دوم تبدیل شود. در اینجا، \en{r2} صرفاً شماره‌ی خطی در فایل دوم را نشان می‌دهد که موقعیت متناظر (یعنی جایی که این سطرها در صورت وجود، در فایل دوم قرار می‌داشتند) را مشخص می‌کند؛ به بیان دیگر، \en{r2} خودِ محتوای حذف‌شده نیست، بلکه صرفاً یک نقطه‌ی مرجع برای مکان‌یابی در فایل دوم است. \\
\bottomrule
\end{tabularx}
\end{table}

در این ساختار، «محدوده» (\en{Range}) شامل شماره‌ی سطر شروع و پایان است که با کاما از یکدیگر جدا می‌شوند. اگرچه این قالب، موسوم به قالب «عادی» (\en{Normal Format})، به دلیل انطباق با استاندارد \en{POSIX} و حفظ سازگاری با نسخه‌های سنتی یونیکس به‌عنوان پیش‌فرض ابزار \cmd{diff} در نظر گرفته شده است، اما در عمل خوانایی چندانی ندارد؛ چرا که هیچ سطر مشترکی از متن اصلی را برای مرجع‌دهی نمایش نمی‌دهد و صرفاً به شماره‌ی خطوط بسنده می‌کند. به همین دلیل، امروزه دو قالب کاربردی‌تر، یعنی قالب «زمینه» (\en{Context Format}) و قالب «یکپارچه» (\en{Unified Format})، رواج بیشتری یافته‌اند.

هنگامی که ابزار \cmd{diff} با قالب «زمینه» (گزینه‌ی \cmd{-c}) اجرا می‌شود، خروجی به شرح زیر نمایش می‌یابد:

\begin{latin}
\begin{lstlisting}
[me@linuxbox ~]$ diff -c file1.txt file2.txt
*** file1.txt	2025-12-23 06:40:13.000000000 -0500
--- file2.txt	2025-12-23 06:40:34.000000000 -0500
***************
*** 1,4 ****
- a
  b
  c
  d
--- 1,4 ----
  b
  c
  d
+ e
\end{lstlisting}
\end{latin}

خروجی با درج نام فایل‌ها و برچسب زمانیِ آخرین تغییرات آن‌ها آغاز می‌گردد. در این ساختار، فایل نخست با نماد ستاره (\cmd{***}) و فایل دوم با نماد خط تیره (\cmd{---}) متمایز می‌شوند. در سراسر گزارش، از همین نمادها برای ارجاع به هر فایل استفاده شده است. پس از بخش سرآیند، گروه‌های تغییرات به همراه تعدادی از سطرهای زمینه (\en{Context Lines}) برای درک بهتر موقعیت تغییر نمایش داده می‌شوند. در اولین گروه، عبارت زیر مشاهده می‌شود:

\begin{latin}
\begin{lstlisting}
*** 1,4 ***
\end{lstlisting}
\end{latin}

این بخش نشان‌دهنده‌ی بازه‌ی سطرهای ۱ تا ۴ در فایل اول است. در ادامه نیز عبارت زیر را مشاهده می‌کنیم:

\begin{latin}
\begin{lstlisting}
--- 1,4 ---
\end{lstlisting}
\end{latin}

این سطر به بازه‌ی متناظر در فایل دوم اشاره دارد. در این قالب، تغییرات با نشانه‌های خاصی در ستون اول مشخص می‌شوند؛ علامت \cmd{-} به معنای حذف سطر از فایل اول، علامت \cmd{+} به معنای افزودن سطر به فایل دوم و علامت \cmd{!} نشان‌دهنده‌ی اصلاح یا تغییر محتوای سطر است.

در قالب «زمینه» (\en{Context Format})، هر سطر درون یک گروه تغییر با یکی از چهار شاخص توصیف‌شده در جدول زیر آغاز می‌شود که وضعیت آن سطر را نسبت به فایل مقابل مشخص می‌کند:

\begin{table}[htbp]
\centering
\caption{شاخص‌های وضعیت سطر در قالب \en{Context}}
\label{tab:context-markers}
\begin{tabularx}{\textwidth}{c X}
\toprule
\textbf{شاخص} & \textbf{معنا و مفهوم فنی} \\
\midrule
(فاصله خالی) & سطر زمینه (\en{Context}): این سطر در هر دو فایل یکسان است و صرفاً برای درک موقعیت تغییر و ایجاد بافت اطلاعاتی (\en{Context}) در اطراف تفاوت‌ها نمایش داده می‌شود. \\
\addlinespace
\cmd{-} & سطر حذفی: این سطر در فایل اول وجود دارد، اما در فایل دوم یافت نمی‌شود و نشان‌دهندهٔ یک خط حذف‌شده است. \\
\addlinespace
\cmd{+} & سطر الحاقی: این سطر در فایل اول وجود ندارد و به فایل دوم افزوده شده است؛ بنابراین یک خط اضافه‌شده محسوب می‌شود. \\
\addlinespace
\cmd{!} & سطر اصلاح‌شده: نشان‌دهندهٔ تغییری است که در یک یا چند سطر رخ داده و محتوای آن‌ها بین دو فایل متفاوت است. در این حالت، هر دو نسخه از سطر (قدیم و جدید) در بخش مربوط به خود در خروجی نمایش داده می‌شوند. \\
\bottomrule
\end{tabularx}
\end{table}

قالب «یکپارچه» (\en{Unified Format}) شباهت ساختاری زیادی به قالب \en{Context} دارد، اما با حذف تکرار سطرهای مشترک، خروجی فشرده‌تر و خواناتری ارائه می‌دهد. این قالب که با گزینه‌ی \cmd{-u} فراخوانی می‌شود، در پروژه‌های نرم‌افزاری مدرن استاندارد رایج‌تری است:

\begin{latin}
\begin{lstlisting}
[me@linuxbox ~]$ diff -u file1.txt file2.txt
--- file1.txt2008-12-23 06:40:13.000000000 -0500
+++ file2.txt2008-12-23 06:40:34.000000000 -0500
@@ -1,4 +1,4 @@
-a
 b
 c
 d
+e
\end{lstlisting}
\end{latin}

تمایز بنیادین میان قالب «زمینه» (\en{Context}) و قالب «یکپارچه» (\en{Unified})، حذف سطرهای تکراری در زمینه (\en{Context}) است؛ امری که موجب می‌گردد خروجی در قالب یکپارچه (\en{Unified}) به‌مراتب کوتاه‌تر و بهینه‌تر باشد. در مثال پیشین، پس از درج برچسب‌های زمانی (مشابه قالب \en{Context})، عبارت زیر درج شده است:

\begin{latin}
\begin{lstlisting}
@@ -1,4 +1,4 @@
\end{lstlisting}
\end{latin}

این رشته مختصات دقیق سطرهایی از فایل اول و دوم را که در این گروهِ تغییر (\en{Hunk}) گنجانده شده‌اند، تبیین می‌کند. در ادامه، خودِ سطرها به همراه سه سطر زمینه (\en{Context Lines}) که به‌صورت پیش‌فرض در اطراف هر گروه تغییر نمایش داده می‌شوند (در صورت وجود)، آورده می‌شوند. هر سطر در این ساختار با یکی از سه نویسه‌ی معرفی‌شده در جدول زیر آغاز می‌گردد:

\begin{table}[htbp]
\centering
\caption{شاخص‌های وضعیت سطر در قالب \en{Unified}}
\label{tab:unified-markers}
\begin{tabularx}{\textwidth}{c X}
\toprule
\textbf{نویسه} & \textbf{مفهوم فنی} \\
\midrule
(فاصله خالی) & سطر مشترک: این سطر در هر دو فایل عیناً تکرار شده و نقش زمینه را دارد. \\
\addlinespace
\cmd{-} & سطر حذفی: این سطر از فایل اول حذف شده است. \\
\addlinespace
\cmd{+} & سطر الحاقی: این سطر به فایل اول افزوده شده تا به محتوای فایل دوم تبدیل شود. \\
\bottomrule
\end{tabularx}
\end{table}

در این ساختار، تمرکز اصلی بر نمایش فشرده‌ی تغییرات است؛ به‌طوری که تحلیل‌گر یا برنامه‌ی پردازش‌گر بتواند به‌سرعت تفاوت‌ها را بدون نیاز به بازخوانیِ مکرر بخش‌های مشترک شناسایی کند.

\subsection{دستور patch}

برنامه‌ی \cmd{patch} به‌منظور اعمال تغییرات روی فایل‌های متنی طراحی شده است. این ابزار خروجی حاصل از \cmd{diff} را به‌عنوان ورودی دریافت کرده و معمولاً برای به‌روزرسانی نسخه‌های قدیمی فایل‌ها به نسخه‌های جدیدتر به‌کار می‌رود. یکی از برجسته‌ترین نمونه‌های کاربرد این ابزار، فرآیند توسعه‌ی هسته‌ی لینوکس (\en{Linux Kernel}) است. هسته‌ی لینوکس توسط شبکه‌ی گسترده‌ای از مشارکت‌کنندگان در سراسر جهان توسعه می‌یابد که به‌طور مستمر اصلاحات و بهبودهای کوچکی را روی کد منبع اعمال می‌کنند.

با توجه به حجم عظیم کد منبع هسته که بالغ بر میلیون‌ها سطر است، ارسال کل ساختار درختی کد منبع (\en{Source Tree}) برای هر تغییر جزئی، رویکردی غیرمنطقی و ناکارآمد است. در عوض، توسعه‌دهندگان یک فایل \cmd{diff} ارسال می‌کنند که تنها شامل تفاوت‌های نسخه‌ی پیشین با نسخه‌ی اصلاح‌شده است. دریافت‌کننده سپس با استفاده از ابزار \cmd{patch} این تغییرات را بر ساختار درختی کد منبع خود اعمال می‌کند. بهره‌گیری از ترکیب \cmd{diff/patch} دو مزیت استراتژیک دارد:

\begin{enumerate}
  \item حجم فایل \cmd{diff} در مقایسه با کل درخت کد منبع بسیار ناچیز است.
  \item فایل \cmd{diff} تغییرات را به‌صورت متمرکز و شفاف نمایش می‌دهد، که این امر فرآیند بازبینی کد (\en{Code Review}) را برای ناظران تسریع می‌کند.
\end{enumerate}

اگرچه این ابزارها در دنیای برنامه‌نویسی به شهرت رسیده‌اند، اما بر روی هر نوع فایل متنی، از جمله فایل‌های پیکربندی سیستم، قابل استفاده هستند. طبق مستندات پروژه \en{GNU}، استاندارد پیشنهادی برای ایجاد یک فایل وصله جهت استفاده در \cmd{patch} به شرح زیر است:

\begin{latin}
\begin{lstlisting}
diff -Naur old_file new_file > diff_file
\end{lstlisting}
\end{latin}

در این دستور، آرگومان‌های ورودی می‌توانند فایل‌های منفرد یا دایرکتوری‌ها باشند. گزینه‌ی \cmd{-r} امکان پیمایش بازگشتی در ساختار دایرکتوری‌ها را فراهم می‌کند. پس از تولید فایل، می‌توان با دستور زیر آن را روی فایل قدیمی اعمال کرد:

\begin{latin}
\begin{lstlisting}
patch < diff_file
\end{lstlisting}
\end{latin}

برای نمایش عملی این فرآیند، از فایل‌های آزمایشی خود استفاده می‌کنیم:

\begin{latin}
\begin{lstlisting}
[me@linuxbox ~]$ diff -Naur file1.txt file2.txt > patchfile.txt
[me@linuxbox ~]$ patch < patchfile.txt
patching file file1.txt
[me@linuxbox ~]$ cat file1.txt
b
c
d
e
\end{lstlisting}
\end{latin}

در این مثال، ابتدا فایل \file{patchfile.txt} را ایجاد و سپس آن را اعمال کردیم. نکته‌ی مهم اینجاست که در هنگام استفاده از \cmd{patch} نیازی به تعیین فایل هدف نیست؛ زیرا فایل \cmd{diff} (در قالب \en{Unified}) نام فایل‌های مربوطه را در بخش سرآیند (\en{Header}) خود شامل می‌شود. همان‌طور که ملاحظه می‌کنید، پس از اتمام عملیات، محتوای \file{file1.txt} دقیقاً با \file{file2.txt} یکسان شده است.

ابزار \cmd{patch} دارای گزینه‌های پیشرفته‌ی متعددی است و برنامه‌های کمکی مکمل دیگری نیز برای تحلیل و ویرایش وصله‌ها وجود دارند.

\section{ویرایش غیرتعاملی متن (Stream Editing)}

تجربه‌ی ما در کار با ویرایشگرهای متن تا به این‌جا عمدتاً به‌صورت تعاملی (\en{Interactive}) بوده است؛ به این معنا که ویرایشگر را باز کرده، مکان‌نما (\en{Cursor}) را به‌صورت دستی جابه‌جا کرده و تغییرات را تایپ نموده‌ایم. با این حال، در دنیای لینوکس و یونیکس، روش‌های غیرتعاملی (\en{Non-interactive}) برای ویرایش متن جایگاه ویژه‌ای دارند.

ویرایش غیرتعاملی زمانی اهمیت می‌یابد که بخواهیم مجموعه‌ای از اصلاحات را بدون دخالت مستقیم کاربر و تنها با اجرای یک فرمان، بر روی یک یا چندین فایل به‌صورت انبوه اعمال کنیم. این رویکرد زیربنای خودکارسازی (\en{Automation}) در مدیریت سیستم است.

\subsection{دستور tr}

برنامه‌ی \cmd{tr} جهت «نویسه‌گردانی» (\en{Transliterate}) یا تبدیل کاراکترها به‌کار می‌رود. این ابزار را می‌توان به‌نوعی عملیاتِ «جست‌وجو و جایگزینی بر پایه‌ی نویسه» در نظر گرفت. نویسه‌گردانی در واقع فرآیند نگاشت کاراکترها از یک مجموعه به مجموعه‌ای دیگر است؛ برای نمونه، تبدیل حروف کوچک به حروف بزرگ (\en{Uppercase}) مثالی کلاسیک از این فرآیند است که با استفاده از \cmd{tr} به شیوه‌ی زیر پیاده‌سازی می‌شود:

\begin{latin}
\begin{lstlisting}
[me@linuxbox ~]$ echo "lowercase letters" | tr a-z A-Z
LOWERCASE LETTERS
\end{lstlisting}
\end{latin}

همان‌طور که در مثال فوق مشهود است، \cmd{tr} منحصراً بر روی ورودی استاندارد (\en{stdin}) عمل کرده و نتایج را در خروجی استاندارد (\en{stdout}) چاپ می‌کند. این برنامه دو آرگومان اصلی می‌پذیرد: نخست، مجموعه‌ای از کاراکترها که باید شناسایی شوند و دوم، مجموعه‌ی متناظری که جایگزین آن‌ها خواهند شد. تعریف مجموعه‌های کاراکتری (\en{Character Sets}) معمولاً به یکی از روش‌های زیر انجام می‌پذیرد:

\begin{itemize}
  \item \textbf{فهرست صریح:} درج مستقیم و تک‌تک کاراکترها؛ برای مثال:
\begin{latin}
\begin{lstlisting}
ABCDEFGHIJKLMNOPQRSTUVWXYZ
\end{lstlisting}
\end{latin}
  \item \textbf{بازه کاراکتری (\en{Range}):} استفاده از ساختارهایی نظیر \cmd{A-Z}. شایان ذکر است که این روش ممکن است تحت تأثیر تنظیمات محلی سیستم (\en{Locale Collation Order}) قرار گیرد، بنابراین در محیط‌های حساس باید با دقت و احتیاط استفاده شود.
  \item \textbf{کلاس‌های نویسه‌ای \en{POSIX}:} روش سوم و استاندارد برای تعریف مجموعه‌ها، استفاده از «کلاس‌های نویسه‌ای» نظیر \cmd{[:upper:]} یا \cmd{[:lower:]} است.
\end{itemize}

در اکثر عملیات‌های نگاشت، طول هر دو مجموعه‌ی کاراکتر باید برابر باشد؛ با این حال، در شرایط خاصی ممکن است مجموعه‌ی اول بزرگ‌تر از دومی تعریف شود؛ به‌ویژه زمانی که قصد داریم چندین کاراکتر مختلف را به یک نویسه‌ی واحد تبدیل کنیم:

\begin{latin}
\begin{lstlisting}
[me@linuxbox ~]$ echo "lowercase letters" | tr [:lower:] A
AAAAAAAAA AAAAAAA
\end{lstlisting}
\end{latin}

علاوه بر قابلیت تبدیل، ابزار \cmd{tr} توانایی حذف نویسه‌های خاص از جریان ورودی را نیز داراست. پیش‌تر در این فصل، به چالش سازگاری فایل‌های متنی سیستم‌عامل \en{MS-DOS} با محیط یونیکس اشاره شد. برای تبدیل این فایل‌ها به فرمت استاندارد یونیکس، ضروری است کاراکتر بازگشت به ابتدای سطر (\en{Carriage Return}) از انتهای هر خط حذف گردد. این عملیات با استفاده از گزینه \cmd{-d} (مخفف \en{Delete}) به سادگی امکان‌پذیر است:

\begin{latin}
\begin{lstlisting}
tr -d '\r' < dos_file > unix_file
\end{lstlisting}
\end{latin}

در این فرمان، \file{dos\_file} به عنوان ورودی و \file{unix\_file} به عنوان خروجیِ اصلاح‌شده در نظر گرفته شده است. توالی گریز \cmd{\textbackslash r} در اینجا نماینده‌ی نویسه‌ی بازگشت به ابتدای سطر (\en{Carriage Return}) است. برای دسترسی به فهرست کامل توالی‌های گریز و کلاس‌های نویسه‌ای پشتیبانی‌شده، می‌توانید از دستور راهنما استفاده کنید:

\begin{latin}
\begin{lstlisting}
[me@linuxbox ~]$ tr --help
\end{lstlisting}
\end{latin}

\subsection{الگوریتم ROT13: پنهان‌سازی متن}

یکی از کاربردهای تفننی و جالب ابزار \cmd{tr} پیاده‌سازی الگوریتم \en{ROT13} است. \en{ROT13} یک روش رمزنگاری ابتدایی مبتنی بر «رمز جانشینی» (\en{Substitution Cipher}) است. در واقع، اطلاق عنوان «رمزنگاری» به \en{ROT13} چندان دقیق نیست و عبارت «پنهان‌سازی متن» (\en{Text Obfuscation}) توصیف واقع‌بینانه‌تری از عملکرد آن ارائه می‌دهد. این روش که گاه برای مخفی کردن محتوای حساس یا پاسخ معماها در تالارهای گفتگو به کار می‌رود، هر حرف انگلیسی را با سیزدهمین حرف پس از آن در الفبا جایگزین می‌کند.

از آنجا که عدد ۱۳ دقیقاً نیمی از حروف الفبای ۲۶گانه‌ی انگلیسی است، این الگوریتم خاصیت «خود-معکوسی» دارد؛ یعنی اعمال مجدد آن بر روی متن رمزگذاری شده، محتوای اصلی را بازیابی می‌کند. برای اجرای این فرآیند با استفاده از \cmd{tr} از دستور زیر استفاده می‌شود:

\begin{latin}
\begin{lstlisting}
[me@linuxbox ~]$ echo "secret text" | tr a-zA-Z n-za-mN-ZA-M
frperg grkg
\end{lstlisting}
\end{latin}

با اجرای مجدد همان دستور بر روی خروجی، متن به حالت اولیه بازمی‌گردد:

\begin{latin}
\begin{lstlisting}
[me@linuxbox ~]$ echo "frperg grkg" | tr a-zA-Z n-za-mN-ZA-M
secret text
\end{lstlisting}
\end{latin}

بسیاری از کلاینت‌های قدیمی ایمیل و خبرخوان‌های یوزنت (\en{Usenet}) به‌صورت پیش‌فرض از این قابلیت پشتیبانی می‌کنند. برای مطالعه‌ی جزئیات بیشتر فنی، مقاله‌ی ویکی‌پدیا در این زمینه منبع ارزشمندی است:

\begin{latin}
\url{http://en.wikipedia.org/wiki/ROT13}
\end{latin}

ابزار \cmd{tr} قابلیت کاربردی دیگری نیز دارد؛ با استفاده از گزینه \cmd{-s} (مخفف \en{Squeeze})، این برنامه می‌تواند تکرارهای متوالی یک نویسه را فشرده کرده و به یک نمونه تقلیل دهد.

\begin{latin}
\begin{lstlisting}
[me@linuxbox ~]$ echo "aaabbbccc" | tr -s ab
abccc
\end{lstlisting}
\end{latin}

در این مثال، مجموعه‌ای از نویسه‌های تکراری به \cmd{tr} داده شده است. با تعیین مجموعه‌ی هدف (در اینجا \cmd{"ab"})، تکرارهای متوالی این حروف حذف شده و تنها یک نسخه از هر کدام باقی می‌ماند، در حالی که نویسه‌ی \cmd{"c"} که در مجموعه‌ی هدف نبود، بدون تغییر باقی مانده است. توجه داشته باشید که عملیات فشرده‌سازی تنها زمانی عمل می‌کند که نویسه‌های تکراری به‌صورت مجاور و پشت‌سرهم باشند. در غیر این صورت، دستور تاثیری بر متن نخواهد داشت:

\begin{latin}
\begin{lstlisting}
[me@linuxbox ~]$ echo "abcabcabc" | tr -s ab
abcabcabc
\end{lstlisting}
\end{latin}

\subsection{ویرایشگر جریانی sed}

نام \cmd{sed} از عبارت \en{Stream Editor} مشتق شده است. این ابزار قدرتمند عملیات ویرایش را مستقیماً بر روی یک جریان متنی (\en{Stream}) انجام می‌دهد که این جریان می‌تواند شامل مجموعه‌ای از فایل‌های مشخص یا داده‌های دریافتی از ورودی استاندارد (\en{stdin}) باشد. \cmd{sed} ابزاری بسیار منعطف و در عین حال پیچیده است (به‌طوری که منابع و کتاب‌های جامعی منحصراً در تبیین قابلیت‌های آن نگاشته شده است)؛ از این رو در این بخش، تنها به بررسی اصول زیربنایی آن بسنده می‌کنیم.

به‌طور کلی، نحوه‌ی عملکرد \cmd{sed} بدین صورت است که یا یک دستور ویرایشی واحد را مستقیماً از خط فرمان دریافت می‌کند و یا آدرس یک «اسکریپت» حاوی مجموعه‌ای از دستورات به آن ارجاع داده می‌شود. سپس \cmd{sed} این فرامین را به‌صورت سطر به سطر بر روی تمامی داده‌های جریان متنی اعمال می‌کند.

در ادامه، نمونه‌ای ساده از عملکرد این ابزار ارائه شده است:

\begin{latin}
\begin{lstlisting}
[me@linuxbox ~]$ echo "front" | sed 's/front/back/'
back
\end{lstlisting}
\end{latin}

در این مثال، ابتدا با استفاده از دستور \cmd{echo} یک جریان متنی تک‌کلمه‌ای ایجاد و آن را به سوی \cmd{sed} هدایت (\en{Pipe}) کرده‌ایم. سپس \cmd{sed} دستور \cmd{s/front/back/} را بر روی متن ورودی اجرا کرده و خروجی \cmd{back} را تولید می‌کند. این دستور کاملاً مشابه فرمان جایگزینی (\en{Substitution}) در ویرایشگر \cmd{vi} عمل می‌کند.

فرامین در \cmd{sed} معمولاً با یک تک‌نویسه آغاز می‌شوند. در مثال پیشین، دستور جایگزینی با نویسه‌ی \cmd{s} مشخص شده است که متعاقب آن، رشته‌های «جست‌وجو» و «جایگزین» قرار می‌گیرند. این رشته‌ها به‌طور سنتی با استفاده از نویسه‌ی اسلش (\cmd{/}) از یکدیگر تفکیک می‌شوند. با این حال، انتخاب جداکننده (\en{Delimiter}) کاملاً اختیاری است؛ \cmd{sed} نخستین نویسه‌ای را که بلافاصله پس از نام فرمان ظاهر شود، به عنوان جداکننده در نظر می‌گیرد. بر همین اساس، می‌توان فرمان قبلی را به شکل زیر نیز بازنویسی کرد:

\begin{latin}
\begin{lstlisting}
[me@linuxbox ~]$ echo "front" | sed 's_front_back_'
back
\end{lstlisting}
\end{latin}

در این نمونه، با قرار دادن نویسه‌ی زیرخط (\cmd{\_}) بلافاصله پس از فرمان، این نویسه نقش جداکننده را ایفا می‌کند. این انعطاف‌پذیری در انتخاب جداکننده، به‌ویژه هنگام کار با مسیرهای فایل که حاوی تعداد زیادی اسلش هستند، به خوانایی و وضوح دستورات کمک شایانی می‌کند.

بسیاری از فرامین \cmd{sed} قابلیت پذیرش یک «نشانی» (\en{Address}) را دارند. نشانی تعیین می‌کند که فرمان ویرایشی دقیقاً بر روی کدام سطر یا سطرها از جریان متنی اعمال شود. در صورت عدم تعیین نشانی، فرمان به‌صورت پیش‌فرض بر تمامی سطرهای جریان اعمال می‌گردد. ساده‌ترین فرم نشانی‌دهی، استفاده از شماره‌ی سطر است. برای درک بهتر، یک نشانی به مثال قبلی اضافه می‌کنیم:

\begin{latin}
\begin{lstlisting}
[me@linuxbox ~]$ echo "front" | sed '1s/front/back/'
back
\end{lstlisting}
\end{latin}

در اینجا، افزودن عدد \cmd{1} به ابتدای فرمان، صراحتاً بیان می‌کند که عملیات جایگزینی تنها باید بر روی سطر نخست اعمال شود. چنانچه نشانی را به عدد دیگری، مثلاً \cmd{2} تغییر دهیم، شاهد خواهیم بود که هیچ تغییری در خروجی رخ نمی‌دهد؛ چرا که ورودی ما در این مثال تنها شامل یک سطر است:

\begin{latin}
\begin{lstlisting}
[me@linuxbox ~]$ echo "front" | sed '2s/front/back/'
front
\end{lstlisting}
\end{latin}

نحوه تعیین نشانی در \cmd{sed} از تنوع بالایی برخوردار است که امکان کنترل دقیق بر فرآیند ویرایش را فراهم می‌کند. جدول زیر متداول‌ترین روش‌های نشانی‌دهی (\en{Addressing}) را تشریح می‌کند:

\begin{table}[htbp]
\centering
\caption{روش‌های نشانی‌دهی در \cmd{sed}}
\label{tab:sed-addressing}
\begin{tabularx}{\textwidth}{l X}
\toprule
\textbf{نشانی} & \textbf{شرح عملکرد فنی} \\
\midrule
\cmd{n} & ارجاع به یک سطر مشخص بر اساس شماره‌ی آن، که در آن \en{n} یک عدد صحیح مثبت است؛ برای نمونه، \cmd{3} به سطر سوم فایل اشاره دارد. \\
\addlinespace
\cmd{\$} & ارجاع به آخرین سطر از جریان متنی (\en{Input Stream}) ورودی. \\
\addlinespace
\cmd{/regexp/} & تطبیق با تمامی سطرهایی که با یک «عبارت باقاعده» (\en{Regular Expression}) مطابقت دارند. عبارت باقاعده به‌صورت پیش‌فرض میان دو نویسه‌ی اسلش (\cmd{/}) قرار می‌گیرد. همچنین می‌توان با استفاده از ساختار \cmd{\textbackslash cregexpc} از یک نویسه‌ی جداکننده‌ی دلخواه به‌جای اسلش بهره برد (که در آن \en{c} نویسه‌ی جداکننده‌ی جایگزین است)؛ این قابلیت به‌ویژه هنگام جست‌وجوی الگوهایی که خود حاوی نویسه‌ی اسلش هستند، مفید است. \\
\addlinespace
\cmd{addr1,addr2} & تعیین یک بازه (\en{Range}) از سطرها، از \en{addr1} تا \en{addr2}، شامل هر دو سطرِ ابتدا و انتها. هر یک از این دو نشانی می‌تواند به‌صورت شماره‌ی سطر، عبارت باقاعده، یا نماد \cmd{\$} بیان شود. \\
\addlinespace
\cmd{first\textasciitilde step} & تطبیق با سطر شماره‌ی \en{first} و سپس تکرار الگو با گام \en{step} (این ساختار یکی از قابلیت‌های افزوده‌ی \en{GNU sed} است). برای نمونه، \cmd{1\textasciitilde 2} تمامی سطرهای فرد را هدف قرار می‌دهد، در حالی که \cmd{5\textasciitilde 5} سطر پنجم و کلیه‌ی مضارب آن (سطرهای ۵، ۱۰، ۱۵ و به همین ترتیب) را انتخاب می‌کند. \\
\addlinespace
\cmd{addr1,+n} & تطبیق با سطر \en{addr1} و \en{n} سطر بلافاصله پس از آن؛ به بیان دیگر، مجموعاً \cmd{n+1} سطر را دربرمی‌گیرد. \\
\addlinespace
\cmd{addr!} & نقیض نشانی (\en{Negation})؛ به معنای اعمال دستور بر تمامی سطرهای فایل، به‌جز آن سطر یا سطرهایی که در \en{addr} مشخص شده‌اند. \\
\bottomrule
\end{tabularx}
\end{table}

این قابلیت‌های نشانی‌دهی، \cmd{sed} را به ابزاری بی‌بدیل برای پردازش‌های دسته‌ای تبدیل می‌کند؛ چرا که اجازه می‌دهد عملیات ویرایش تنها بر بخش‌های ساختاریافته‌ی فایل اعمال شود، بدون آنکه نیازی به پیمایش دستی متن باشد.

در ادامه، نحوه‌ی به‌کارگیری نشانی‌ها در سناریوهای مختلف را بررسی می‌کنیم. ابتدا، تعیین یک بازه‌ی مشخص از شماره‌ی سطرها را در نظر بگیرید:

\begin{latin}
\begin{lstlisting}
[me@linuxbox ~]$ sed -n '1,5p' distros.txt
SUSE     10.2    12/07/2006
Fedora   10      11/25/2008
SUSE     11.0    06/19/2008
Ubuntu   8.04    04/24/2008
Fedora   8       11/08/2007
\end{lstlisting}
\end{latin}

در این مثال، سطرهای ۱ تا ۵ فایل چاپ شده‌اند. برای دستیابی به این نتیجه، از دستور \cmd{p} (مخفف \en{Print}) جهت چاپ سطرهای منطبق بر نشانی استفاده کردیم. نکته‌ی کلیدی در اینجا استفاده از گزینه‌ی \cmd{-n} (حالت \en{Suppress} یا عدم چاپ خودکار) است؛ این گزینه مانع از آن می‌شود که \cmd{sed} طبق رفتار پیش‌فرض خود، تمامی سطرهای ورودی را در خروجی تکرار کند.

سپس، نشانی‌دهی بر پایه‌ی «عبارت باقاعده» (\en{Regular Expression}) را بررسی می‌کنیم:

\begin{latin}
\begin{lstlisting}
[me@linuxbox ~]$ sed -n '/SUSE/p' distros.txt
SUSE     10.2    12/07/2006
SUSE     11.0    06/19/2008
SUSE     10.3    10/04/2007
SUSE     10.1    05/11/2006
\end{lstlisting}
\end{latin}

با محصور کردن الگوی \cmd{/SUSE/} میان دو اسلش، توانستیم تنها سطرهایی را که حاوی این عبارت هستند استخراج کنیم؛ عملکردی که شباهت بسیاری به ابزار \cmd{grep} دارد.

در نهایت، با استفاده از علامت تعجب (\cmd{!})، عملیات «نقیض‌سازی» یا منفی‌سازی نشانی (\en{Address Inversion}) را اعمال می‌کنیم:

\begin{latin}
\begin{lstlisting}
[me@linuxbox ~]$ sed -n '/SUSE/!p' distros.txt
Fedora   10      11/25/2008
Ubuntu   8.04    04/24/2008
Fedora   8       11/08/2007
Ubuntu   6.10    10/26/2006
Fedora   7       05/31/2007
Ubuntu   7.10    10/18/2007
Ubuntu   7.04    04/19/2007
Fedora   6       10/24/2006
Fedora   9       05/13/2008
Ubuntu   6.06    06/01/2006
Ubuntu   8.10    10/30/2008
Fedora   5       03/20/2006
\end{lstlisting}
\end{latin}

خروجی فوق دقیقاً همان چیزی است که انتظار می‌رفت: تمامی سطرهای فایل نمایش داده شده‌اند، به‌استثنای مواردی که با عبارت باقاعده مشخص‌شده مطابقت داشتند.

تا کنون با دو فرمان محوری \cmd{s} و \cmd{p} در \cmd{sed} آشنا شده‌ایم. جدول زیر مجموعه‌ای جامع‌تر از دستورات ویرایشی پایه در این ابزار را معرفی می‌کند:

\begin{table}[htbp]
\centering
\caption{دستورات ویرایشی پایه در \cmd{sed}}
\label{tab:sed-commands}
\begin{tabularx}{\textwidth}{l X}
\toprule
\textbf{دستور} & \textbf{شرح عملکرد فنی} \\
\midrule
\cmd{=} & چاپ شماره‌ی سطر جاری در خروجی استاندارد، پیش از چاپ خودِ محتوای آن سطر. \\
\addlinespace
\cmd{a} & الحاق (\en{Append}) یک متن مشخص، بلافاصله پس از سطر جاری، به جریان خروجی. \\
\addlinespace
\cmd{d} & حذف (\en{Delete}) سطر جاری از جریان خروجی؛ سطر حذف‌شده در نتیجه‌ی نهایی نمایش داده نمی‌شود. \\
\addlinespace
\cmd{i} & درج (\en{Insert}) یک متن مشخص، بلافاصله پیش از سطر جاری، در جریان خروجی. \\
\addlinespace
\cmd{p} & چاپ (\en{Print}) سطر جاری. از آنجا که \cmd{sed} به‌طور پیش‌فرض تمامی سطرها را در خروجی چاپ می‌کند، این دستور معمولاً همراه با گزینه‌ی \cmd{-n} (غیرفعال‌سازی چاپ خودکار) به‌کار می‌رود تا کنترل دقیقی بر خروجی نهایی اعمال شود. \\
\addlinespace
\cmd{q} & خروج (\en{Quit}) از برنامه و توقف پردازش سطرهای باقی‌مانده. در صورتی که گزینه‌ی \cmd{-n} فعال نباشد، سطر جاری پیش از خروج چاپ می‌شود. \\
\addlinespace
\cmd{Q} & خروج فوری از برنامه، بدون پردازش سطرهای باقی‌مانده و بدون چاپ سطر جاری، حتی در حالت پیش‌فرض. \\
\addlinespace
\cmd{s/regexp/replacement/} & جایگزینی (\en{Substitute}) بخشی از متن هر سطر که با الگوی \file{regexp} مطابقت دارد، با متنِ \file{replacement}. در بخش \file{replacement} می‌توان از نویسه‌ی ویژه‌ی \cmd{\&} برای ارجاع به کل متن تطبیق‌یافته استفاده کرد، و همچنین از ارجاع‌های بازگشتی (\en{Backreferences}) مانند \cmd{\textbackslash 1} تا \cmd{\textbackslash 9} برای فراخوانی محتوای گروه‌های مشخص‌شده در عبارت باقاعده بهره برد. افزودن پرچم‌هایی (\en{Flags}) در انتهای دستور، مانند \cmd{g} برای جایگزینی سراسری در کل سطر، رفتار پیش‌فرض این دستور را تغییر می‌دهد. \\
\addlinespace
\cmd{y/set1/set2} & نویسه‌گردانی (\en{Transliterate})، یعنی تبدیل یک‌به‌یکِ هر نویسه در مجموعه‌ی \file{set1} به نویسه‌ی متناظر آن در مجموعه‌ی \file{set2}. برخلاف ابزار \cmd{tr}، در این دستور الزامی است که طول دو مجموعه‌ی \file{set1} و \file{set2} دقیقاً یکسان باشد. \\
\bottomrule
\end{tabularx}
\end{table}

بدون شک دستور \cmd{s} (جایگزینی) پرکاربردترین فرمان ویرایشی در \cmd{sed} محسوب می‌شود. برای نمایش قدرت این ابزار، تغییری را بر روی فایل \file{distros.txt} اعمال می‌کنیم. همان‌طور که پیش‌تر اشاره شد، قالب تاریخ در این فایل به‌صورت \en{MM/DD/YYYY} درج شده که برای پردازش‌های سیستمی و عملیات مرتب‌سازی مناسب نیست؛ چرا که برای مرتب‌سازی بهینه، قالب استاندارد \en{YYYY-MM-DD} ترجیح داده می‌شود. اصلاح دستی این موارد نه‌تنها زمان‌بر، بلکه مستعد خطای انسانی است، اما با \cmd{sed} می‌توان این تغییر را تنها در یک گام عملیاتی کرد:

\begin{latin}
\begin{lstlisting}
[me@linuxbox ~]$ sed 's/\([0-9]\{2\}\)\/\([0-9]\{2\}\)\/\([0-9]\{4\}\)$/\3-\1-\2/' distros.txt
SUSE     10.2    2006-12-07
Fedora   10      2008-11-25
SUSE     11.0    2008-06-19
Ubuntu   8.04    2008-04-24
Fedora   8       2007-11-08
SUSE     10.3    2007-10-04
Ubuntu   6.10    2006-10-26
Fedora   7       2007-05-31
Ubuntu   7.10    2007-10-18
Ubuntu   7.04    2007-04-19
SUSE     10.1    2006-05-11
Fedora   6       2006-10-24
Fedora   9       2008-05-13
Ubuntu   6.06    2006-06-01
Ubuntu   8.10    2008-10-30
Fedora   5       2006-03-20
\end{lstlisting}
\end{latin}

ظاهر این دستور ممکن است در ابتدا غریب و پیچیده به نظر برسد، اما عملکرد آن دقیق است. این مثال به‌خوبی نشان می‌دهد که چرا برخی به شوخی عبارات باقاعده (\en{Regular Expressions}) را «فقط‌نوشتنی» می‌نامند؛ به این معنا که نوشتن آن‌ها ساده‌تر از بازخوانی و درک مجدد آن‌هاست. پیش از آنکه از پیچیدگی این دستور دلسرد شوید، بیایید اجزای آن را کالبدشکافی کنیم. ساختار بنیادی این فرمان به شرح زیر است:

\begin{latin}
\begin{lstlisting}
sed 's/regexp/replacement/' distros.txt
\end{lstlisting}
\end{latin}

گام نخست، طراحی یک عبارت باقاعده برای شناسایی و تفکیک اجزای تاریخ است. با توجه به اینکه تاریخ در انتهای سطر و به فرمت \en{MM/DD/YYYY} قرار دارد، می‌توانیم از الگویی مشابه زیر استفاده کنیم:

\begin{latin}
\begin{lstlisting}
[0-9]{2}/[0-9]{2}/[0-9]{4}$
\end{lstlisting}
\end{latin}

الگوی فوق به‌طور دقیق دو رقم، یک اسلش، دو رقم، اسلشی دیگر، چهار رقم و در نهایت انتهای سطر را شناسایی و تطبیق می‌دهد. با مشخص شدن بخش \file{regexp} (الگو)، نوبت به طراحی بخش \file{replacement} (جایگزین) می‌رسد. برای این کار، باید از قابلیتی در عبارات باقاعده موسوم به ارجاع‌های بازگشتی (\en{Backreferences}) استفاده کنیم. این ویژگی که در ابزارهای مبتنی بر \en{BRE} در دسترس است، به این صورت عمل می‌کند: چنانچه توالی \cmd{\textbackslash n} (که در آن \en{n} عددی بین ۱ تا ۹ است) در بخش جایگزین ظاهر شود، به زیرعبارت (\en{Sub-expression}) متناظر در بخش الگو ارجاع می‌دهد. برای ایجاد این زیرعبارات، کافی است بخش‌های مورد نظر را درون پرانتز قرار دهیم:

\begin{latin}
\begin{lstlisting}
([0-9]{2})/([0-9]{2})/([0-9]{4})$
\end{lstlisting}
\end{latin}

اکنون سه زیرعبارت مجزا داریم: اولی ماه، دومی روز و سومی سال را در خود جای داده است. با این ساختار، می‌توانیم بخش \file{replacement} را به شکل زیر تدوین کنیم:

\begin{latin}
\begin{lstlisting}
\3-\1-\2
\end{lstlisting}
\end{latin}

این ترکیب، ترتیب قرارگیری را به «روز-ماه-سال» تغییر می‌دهد. با این حال، دستور نهایی ما به شکل زیر با دو چالش فنی روبروست:

\begin{latin}
\begin{lstlisting}
sed 's/([0-9]{2})/([0-9]{2})/([0-9]{4})$/\3-\1-\2/' distros.txt
\end{lstlisting}
\end{latin}

چالش اول، وجود اسلش‌ها (\cmd{/}) در درون تاریخ است که باعث تداخل با جداکننده‌های دستور \cmd{s} شده و \cmd{sed} را در تشخیص مرزهای دستور دچار سردرگمی می‌کند. چالش دوم این است که \cmd{sed} به‌صورت پیش‌فرض از نحو \en{BRE} (\en{Basic Regular Expressions}) پیروی می‌کند؛ در این استاندارد، پرانتزها و آکولادها بدون نویسه‌ی گریز، صرفاً به عنوان کاراکترهای عادی در نظر گرفته می‌شوند و نقش متاکاراکتر (\en{Metacharacter}) ایفا نمی‌کنند.

برای رفع این دو مشکل، از نویسه‌ی بک‌اسلش (\cmd{\textbackslash}) جهت گریزدهی (\en{Escaping}) نویسه‌های خاص استفاده می‌کنیم تا معنای فنی آن‌ها برای مفسر روشن شود:

\begin{latin}
\begin{lstlisting}
sed 's/\([0-9]\{2\}\)\/\([0-9]\{2\}\)\/\([0-9]\{4\}\)$/\3-\1-\2/' distros.txt
\end{lstlisting}
\end{latin}

و بدین ترتیب، عملیات اصلاح ساختار تاریخ با موفقیت به پایان می‌رسد.

قابلیت دیگر دستور \cmd{s} بهره‌گیری از پرچم‌های اختیاری (\en{Flags}) در انتهای رشته‌ی جایگزین است. محوری‌ترینِ این پرچم‌ها، \cmd{g} (مخفف \en{Global}) است که به \cmd{sed} فرمان می‌دهد عملیات جایگزینی را به‌صورت سراسری در کل سطر انجام دهد؛ در حالی که رفتار پیش‌فرض ابزار، تنها اعمال تغییر بر روی نخستین موردِ انطباق یافته است.

به مثال زیر توجه کنید:

\begin{latin}
\begin{lstlisting}
[me@linuxbox ~]$ echo "aaabbbccc" | sed 's/b/B/'
aaaBbbccc
\end{lstlisting}
\end{latin}

مشاهده می‌شود که جایگزینی صرفاً برای اولین کاراکتر \cmd{b} رخ داده و سایر موارد بدون تغییر مانده‌اند. با الحاق پرچم \cmd{g} به انتهای دستور، می‌توان تمامی موارد منطبق را تغییر داد:

\begin{latin}
\begin{lstlisting}
[me@linuxbox ~]$ echo "aaabbbccc" | sed 's/b/B/g'
aaaBBBccc
\end{lstlisting}
\end{latin}

تا به این‌جا، تنها از دستورات تک‌خطی در رابط خط فرمان استفاده کرده‌ایم. با این حال، می‌توان مجموعه‌ای از دستورات پیچیده را در قالب یک اسکریپت (\en{Script File}) سازماندهی و با استفاده از گزینه‌ی \cmd{-f} اجرا نمود. برای نمایش این قابلیت، گزارشی از فایل \file{distros.txt} تهیه می‌کنیم که شامل یک عنوان در ابتدا، اصلاح فرمت تاریخ‌ها و تبدیل نام توزیع‌ها به حروف بزرگ (\en{Uppercase}) باشد. برای این منظور، یک فایل متنی ایجاد کرده و دستورات زیر را در آن درج می‌کنیم:

\begin{latin}
\begin{lstlisting}
# sed script to produce Linux distributions report

1 i\
\
Linux Distributions Report\
s/\([0-9]\{2\}\)\/\([0-9]\{2\}\)\/\([0-9]\{4\}\)$/\3-\1-\2/
y/abcdefghijklmnopqrstuvwxyz/ABCDEFGHIJKLMNOPQRSTUVWXYZ/
\end{lstlisting}
\end{latin}

در این اسکریپت، از دستور \cmd{i} جهت درج عنوان در سطر نخست، از فرمان \cmd{s} برای بازآرایی ساختار تاریخ و از دستور \cmd{y} برای نویسه‌گردانی حروف کوچک به بزرگ استفاده شده است.

پس از ذخیره‌سازی اسکریپت با نام \file{distros.sed} (با پسوند متداول فایل‌های اسکریپت \cmd{sed})، آن را به شکل زیر بر روی فایل هدف اجرا می‌کنیم:

\begin{latin}
\begin{lstlisting}
[me@linuxbox ~]$ sed -f distros.sed distros.txt

Linux Distributions Report

SUSE     10.2    2006-12-07
FEDORA   10      2008-11-25
SUSE     11.0    2008-06-19
UBUNTU   8.04    2008-04-24
FEDORA   8       2007-11-08
SUSE     10.3    2007-10-04
UBUNTU   6.10    2006-10-26
FEDORA   7       2007-05-31
UBUNTU   7.10    2007-10-18
UBUNTU   7.04    2007-04-19
SUSE     10.1    2006-05-11
FEDORA   6       2006-10-24
FEDORA   9       2008-05-13
UBUNTU   6.06    2006-06-01
UBUNTU   8.10    2008-10-30
FEDORA   5       2006-03-20
\end{lstlisting}
\end{latin}

همان‌طور که مشاهده می‌شود، اسکریپت خروجی مطلوب را تولید کرده است؛ اما مکانیسم عملکرد آن چگونه است؟ برای درک بهتر لایه‌های اجرایی، محتوای اسکریپت را با استفاده از دستور \cmd{cat} و گزینه شماره‌گذاری سطرها بازخوانی می‌کنیم:

\begin{latin}
\begin{lstlisting}
[me@linuxbox ~]$ cat -n distros.sed
     1  # sed script to produce Linux distributions report
     2
     3  1 i\
     4  \
     5  Linux Distributions Report\
     6
     7  s/\([0-9]\{2\}\)\/\([0-9]\{2\}\)\/\([0-9]\{4\}\)$/\3-\1-\2/
     8  y/abcdefghijklmnopqrstuvwxyz/ABCDEFGHIJKLMNOPQRSTUVWXYZ/
\end{lstlisting}
\end{latin}

سطر نخست اسکریپت ما یک «توضیح» (\en{Comment}) است. مشابه بسیاری از فایل‌های پیکربندی و زبان‌های برنامه‌نویسی در محیط لینوکس، توضیحات با کاراکتر \cmd{\#} آغاز می‌شوند و متنی خوانا برای انسان را در بر می‌گیرند. توضیحات می‌توانند در هر بخشی از اسکریپت (به‌جز در میانه‌ی یک دستور) قرار گیرند و برای مستندسازی و نگهداری کد در آینده اهمیت بسزایی دارند.

سطر دوم یک سطر خالی است. سطرهای خالی، همانند توضیحات، صرفاً جهت تفکیک منطقی بخش‌های مختلف و افزایش خوانایی اسکریپت به کار می‌روند و هیچ تأثیری بر عملکرد آن ندارند.

بسیاری از دستورات \cmd{sed} قابلیت پذیرش «نشانی» (\en{Address}) را دارند. این نشانی‌ها تعیین می‌کنند که عملیات ویرایشی دقیقاً بر روی کدام بخش از ورودی اعمال شود. نشانی می‌تواند یک شماره سطر واحد، بازه‌ای از سطرها، یا نویسه‌ی ویژه‌ی \cmd{\$} باشد که به آخرین سطر جریان ورودی اشاره دارد.

سطر ۳ اسکریپت، دستور \cmd{i} را با نشانی \cmd{1} (یعنی نخستین سطر ورودی) فراخوانی می‌کند؛ به این ترتیب، متنی که در ادامه می‌آید، پیش از نخستین سطر ورودی درج (\en{Insert}) خواهد شد. این دستور از دنباله‌ای متشکل از بک‌اسلش (\cmd{\textbackslash}) و به‌دنبال آن یک نویسه‌ی خط‌جدید (\en{Newline}) استفاده می‌کند تا یک «ادامه خط» (\en{Line Continuation}) ایجاد کند. این مکانیزم، که در اسکریپت‌نویسی پوسته نیز رایج است، اجازه می‌دهد نویسه‌ی خط‌جدید در دل متن تعبیه شود، بدون آنکه مفسر (در اینجا \cmd{sed}) آن را به معنای پایان دستور تلقی کند. دستورات \cmd{i} (درج)، \cmd{a} (الحاق) و \cmd{c} (جایگزینی سطر) امکان افزودن چندین سطر متن را فراهم می‌کنند؛ به این شرط که تمامی سطرها، به‌جز سطر پایانی، با نویسه‌ی ادامه‌خط (\cmd{\textbackslash}) خاتمه یابند. بر همین اساس، متن درج‌شده در سطرهای ۴ تا ۶ اسکریپت ما قرار دارد: سطر ۴ صرفا از یک نویسه‌ی بک‌اسلش تشکیل شده که با ایفای نقش ادامه‌خط، یک سطر خالی به متنِ درج‌شده می‌افزاید؛ سطر ۵ عبارت \en{``Linux Distributions Report''} را در بر دارد؛ و سطر ۶، به دستور \cmd{i} خاتمه می‌دهد.

\begin{note}
\textbf{نکته:} نویسه‌ی ادامه‌خط تنها زمانی معتبر است که بک‌اسلش دقیقاً و بدون هیچ فاصله‌ای، پیش از خط‌جدید (\en{Newline}) قرار گیرد؛ درج هرگونه فاصله‌ی اضافی (\en{Space}) میان بک‌اسلش و نویسه‌ی خط‌جدید، اعتبار این دنباله را از بین می‌برد و می‌تواند باعث بروز خطا در اجرای اسکریپت شود.
\end{note}

سطر ۷ حاوی دستور جست‌وجو و جایگزینی (\cmd{s}) است. از آنجا که هیچ نشانیِ مشخصی برای این دستور تعیین نشده، تمامی سطرهای جریان ورودی مشمول این عملیات می‌شوند.

سطر ۸ عملیات نویسه‌گردانی، یعنی تبدیل حروف کوچک به بزرگ، را بر عهده دارد. شایان ذکر است که برخلاف ابزار \cmd{tr}، دستور \cmd{y} در \cmd{sed} از بازه‌های نویسه‌ای (نظیر \cmd{[a-z]}) یا کلاس‌های نویسه‌ای \en{POSIX} پشتیبانی نمی‌کند و باید تمامی نویسه‌ها به‌صورت صریح و یک‌به‌یک ذکر شوند. همچنین، به دلیل عدم تعیین نشانی برای این دستور نیز، عملیات نویسه‌گردانی بر تمامی سطرهای جریان ورودی اعمال می‌گردد.

\subsection{علاقه‌مندان به ابزار sed معمولاً به ابزارهای زیر نیز تمایل نشان می‌دهند}

\cmd{sed} ابزاری بسیار توانمند است و قابلیت اجرای ویرایش‌های نسبتاً پیچیده بر روی جریان‌های متنی را دارد. با این حال، کاربرد غالب آن بیشتر در قالب دستورات تک‌خطی (\en{One-liners}) است تا نگارش اسکریپت‌های حجیم و پیچیده. برای پروژه‌های بزرگ‌تر و پردازش‌های سنگین‌تر، بسیاری از کاربران ترجیح می‌دهند به سراغ زبان‌های برنامه‌نویسیِ کامل‌تری بروند که در میان آن‌ها، \cmd{awk} و \cmd{Perl} از محبوبیت بیشتری برخوردارند. این دو، برخلاف \cmd{sed} که صرفاً یک ابزار پردازش جریان متنی (\en{Stream Editor}) است، زبان‌های برنامه‌نویسی مستقلی به شمار می‌روند که امکانات به‌مراتب گسترده‌تری، از جمله متغیرها، ساختارهای کنترلی و توابع، در اختیار کاربر قرار می‌دهند.

زبان \en{Perl} به‌طور گسترده‌ای به‌عنوان جایگزینی برای اسکریپت‌نویسی پوسته در وظایف مدیریت سیستم، و همچنین در حوزه‌ی توسعه‌ی وب، مورد استفاده قرار می‌گیرد. در مقابل، \cmd{awk} که خود نیز یک زبان برنامه‌نویسیِ مستقل محسوب می‌شود، کاربردی تخصصی‌تر دارد و قدرت اصلی‌اش در پردازش و تحلیل داده‌های جدولی (\en{Tabular Data}) نهفته است. \cmd{awk} از جهاتی به \cmd{sed} شباهت دارد؛ از جمله اینکه فایل‌ها را سطر به سطر پردازش می‌کند و از ساختاری مشابه با الگوی «نشانی و سپس عملیات» (\en{Address-Action}) بهره می‌برد.

اگرچه بررسی تخصصیِ \cmd{awk} و \cmd{Perl} در چارچوب مباحث این کتاب نمی‌گنجد، اما تسلط بر آن‌ها برای کاربران حرفه‌ای خط فرمان لینوکس، مهارتی بسیار ارزشمند و راهبردی محسوب می‌شود.

\subsection{بررسی‌گر املایی aspell}

آخرین ابزاری که بررسی خواهیم کرد، \cmd{aspell} است؛ یک غلط‌یاب املایی تعاملی (\en{Interactive Spell Checker}) که به عنوان جانشین مدرن برنامه‌ی قدیمی‌ترِ \cmd{ispell} توسعه یافته است و در اکثر موارد می‌تواند به عنوان جایگزینی مستقیم برای آن عمل کند. هرچند \cmd{aspell} عمدتاً به عنوان موتور زیرساختی در سایر نرم‌افزارها برای افزودن قابلیت غلط‌یابی املایی به کار می‌رود، اما به عنوان یک ابزار مستقل در خط فرمان نیز بسیار کارآمد است. این برنامه از هوشمندی بالایی برخوردار است و می‌تواند انواع مختلفی از اسناد فنی از جمله فایل‌های \en{HTML}، کدهای منبع \en{C/C++}، پیام‌های ایمیل و سایر متون تخصصی را به‌دقت تحلیل کند.

جهت بررسی املایی یک فایل متنی ساده، دستور زیر به کار می‌رود:

\begin{latin}
\begin{lstlisting}
aspell check textfile
\end{lstlisting}
\end{latin}

در این دستور، \file{textfile} نشان‌دهنده نام فایلی است که قصد بازبینی و اصلاح آن را دارید. پس از اجرای این فرمان، \cmd{aspell} محیطی تعاملی را می‌گشاید که در آن کلمات مشکوک برجسته شده و پیشنهادهای اصلاحی ارائه می‌گردد.

به عنوان یک مثال عملی، فایلی متنی به نام \file{foo.txt} ایجاد می‌کنیم که حاوی چند خطای املایی عمدی باشد:

\begin{latin}
\begin{lstlisting}
[me@linuxbox ~]$ cat > foo.txt
The quick brown fox jimps over the laxy dog.
\end{lstlisting}
\end{latin}

سپس فرایند عیب‌یابی را با استفاده از \cmd{aspell} آغاز می‌کنیم:

\begin{latin}
\begin{lstlisting}
[me@linuxbox ~]$ aspell check foo.txt
\end{lstlisting}
\end{latin}

از آنجا که \cmd{aspell} در حالت تعاملی (\en{Interactive Mode}) اجرا می‌شود، رابط کاربری مشابه زیر نمایان خواهد شد:

\begin{latin}
\begin{lstlisting}
The quick brown fox jimps over the laxy dog.

------------------------------------------
1) jumps                  6) limps
2) imps                   7) pimps
3) gimps                  8) wimps
4) jump's                 9) comps
5) Jim's                  0) gimp's
i) Ignore                 I) Ignore all
r) Replace                R) Replace all
a) Add                    l) Add Lower
b) Abort                  x) Exit
------------------------------------------
?
\end{lstlisting}
\end{latin}

در بخش فوقانی این رابط، متن فایل نمایش داده شده و واژه‌ی مشکوک به خطای املایی به‌صورت برجسته (\en{Highlight}) مشخص می‌گردد. در بخش میانی، ده واژه‌ی پیشنهادی با شماره‌های \cmd{0} تا \cmd{9} به همراه فهرستی از عملیات‌های ممکن ارائه شده است. در نهایت، در پایین‌ترین قسمت سطر، محیط اعلان (\en{Prompt}) منتظر انتخاب و ورودی کاربر باقی می‌ماند.

با انتخاب کلید \cmd{1}، ابزار \cmd{aspell} واژه‌ی نادرست را با «\en{jumps}» جایگزین نموده و به‌طور خودکار به سراغ خطای بعدی، یعنی «\en{laxy}» می‌رود. در صورت انتخاب گزینه‌ی «\en{lazy}»، فرآیند اصلاح تکمیل شده و برنامه خاتمه می‌یابد. پس از پایان عملیات، با بازخوانی فایل می‌توان مشاهده کرد که اصلاحات با موفقیت اعمال شده‌اند:

\begin{latin}
\begin{lstlisting}
[me@linuxbox ~]$ cat foo.txt
The quick brown fox jumps over the lazy dog.
\end{lstlisting}
\end{latin}

شایان ذکر است که \cmd{aspell} به‌صورت پیش‌فرض یک نسخه‌ی پشتیبان (\en{Backup}) از فایل اصلی با پسوند \file{.bak} ایجاد می‌کند، مگر آنکه با استفاده از گزینه \cmd{--dont-backup} از این کار جلوگیری شود.

برای بررسی توانمندی‌های ویرایشی \cmd{sed} و بازگرداندن فایل به حالت اولیه جهت آزمایش‌های بعدی، از دستور زیر استفاده می‌کنیم:

\begin{latin}
\begin{lstlisting}
[me@linuxbox ~]$ sed -i 's/lazy/laxy/; s/jumps/jimps/' foo.txt
\end{lstlisting}
\end{latin}

در این دستور، گزینه‌ی \cmd{-i} (مخفف \en{In-place}) به \cmd{sed} فرمان می‌دهد که تغییرات را مستقیماً بر روی فایل اصلی بازنویسی کند، به‌جای آنکه خروجی را تنها در ترمینال نمایش دهد. همچنین ملاحظه می‌کنید که می‌توان چندین فرمان ویرایشی را با استفاده از سمی‌کالن (\cmd{;}) در یک خط ترکیب کرد.

در ادامه، چگونگی مدیریت قالب‌های مختلف متنی توسط \cmd{aspell} را بررسی می‌کنیم. با استفاده از یک ویرایشگر متن (نظیر \cmd{vim}) یا ابزار \cmd{sed} مورد بحث، نشانه‌گذاری‌های \en{HTML} را به فایل خود اضافه می‌کنیم:

\begin{latin}
\begin{lstlisting}
<html>
  <head>
    <title>Mispelled HTML file</title>
  </head>
  <body>
    <p>The quick brown fox jimps over the laxy dog.</p>
  </body>
</html>
\end{lstlisting}
\end{latin}

چنانچه بخواهیم فایل اصلاح‌شده را در حالت عادی بررسی کنیم، با یک چالش فنی مواجه می‌شویم. اجرای دستور به صورت پیش‌فرض:

\begin{latin}
\begin{lstlisting}
[me@linuxbox ~]$ aspell check foo.txt
\end{lstlisting}
\end{latin}

منجر به نمایش خروجی زیر می‌گردد:

\begin{latin}
\begin{lstlisting}
<html>
    <head>
            <title>Mispelled HTML file</title>
    </head>
    <body>
            <p>The quick brown fox jimps over the laxy dog.</p>
    </body>
</html>

-------------------------------------
1) HTML                  4) Hamel
2) ht ml                 5) Hamil
3) ht-ml                 6) hotel
i) Ignore                I) Ignore all
r) Replace               R) Replace all
a) Add                   l) Add Lower
b) Abort                 x) Exit
-------------------------------------
?
\end{lstlisting}
\end{latin}

در این وضعیت، \cmd{aspell} برچسب‌های (\en{Tags}) ساختاری \en{HTML} را به عنوان واژگان متنی تلقی کرده و آن‌ها را اشتباه تشخیص می‌دهد. برای رفع این تداخل، باید از گزینه \cmd{-H} (مخفف \en{HTML}) استفاده کرد تا حالت بررسی هوشمند اسناد وب فعال شود:

\begin{latin}
\begin{lstlisting}
[me@linuxbox ~]$ aspell -H check foo.txt
\end{lstlisting}
\end{latin}

خروجی در این حالت به شکل زیر تغییر می‌یابد:

\begin{latin}
\begin{lstlisting}
<html>
      <head>
              <title>Mispelled HTML file</title>
     </head>
     <body>
              <p>The quick brown fox jimps over the laxy dog.</p>
     </body>

---------------------------------------
</html>
1) Mi spelled              6) Misapplied
2) Mi-spelled              7) Miscalled
3) Misspelled              8) Respelled
4) Dispelled               9) Misspell
5) Spelled                 0) Misled
i) Ignore                  I) Ignore all
r) Replace                 R) Replace all
a) Add                     l) Add Lower
b) Abort                   x) Exit
---------------------------------------
?
\end{lstlisting}
\end{latin}

در این پیکربندی، \cmd{aspell} از نشانه‌گذاری‌های \en{HTML} صرف‌نظر کرده و تمرکز خود را تنها بر محتوای متنیِ سند معطوف می‌کند. شایان ذکر است که این ابزار به قدری هوشمند عمل می‌کند که مفادِ ویژگی‌هایی نظیر \en{ALT} را (که حاوی متن توصیفی هستند) بررسی کرده، اما از تگ‌های ساختاری عبور می‌کند.

\begin{note}
\textbf{نکته:} \cmd{aspell} به‌صورت پیش‌فرض از بررسی آدرس‌های اینترنتی (\en{URLs}) و نشانی‌های ایمیل خودداری می‌کند. این رفتار عملیاتی از طریق گزینه‌های خط فرمان قابل سفارشی‌سازی است؛ همچنین کاربر می‌تواند تگ‌های خاصی را برای بررسی یا نادیده گرفتن استثنا کند. جهت مطالعه‌ی جزئیات دقیق‌تر، می‌توانید به صفحه‌ی راهنمای سیستمی (\cmd{man aspell}) مراجعه کنید.
\end{note}

\section{جمع‌بندی}

در این فصل، به واکاوی مجموعه‌ای از ابزارهای خط فرمان پرداختیم که به‌طور اختصاصی برای پردازش و دستکاری متون طراحی شده‌اند. در فصل آتی، ابزارهای تکمیلی دیگری را نیز در این حوزه معرفی خواهیم کرد. اگرچه ممکن است در نگاه نخست، ضرورت یا نحوه‌ی به‌کارگیری روزمره‌ی برخی از این ابزارها چندان ملموس به نظر نرسد، اما تلاش کردیم با ارائه‌ی نمونه‌های عملی، بخشی از توانمندی‌های آن‌ها را در سناریوهای واقعی به تصویر بکشیم.

در فصول پیش‌رو خواهیم دید که این فرامین، ارکان اصلی یک «جعبه‌ابزار» تخصصی را تشکیل می‌دهند که برای حل طیف گسترده‌ای از چالش‌های عملیاتی در لینوکس به‌کار می‌روند. اهمیت و ارزش راهبردی این ابزارها به‌ویژه هنگام ورود به مبحث اسکریپت‌نویسی پوسته (\en{Shell Scripting}) بیش از پیش نمایان خواهد شد؛ جایی که ترکیب این دستورات ساده، امکان اتوماسیون فرآیندهای پیچیده و مدیریت بهینه‌ی سیستم را فراهم می‌سازد.

\section{منابع مطالعه تکمیلی}

وب‌سایت پروژه گنو (\en{GNU}) میزبان مستندات آنلاین جامع و مرجعی برای ابزارهای بررسی‌شده در این فصل است:

از بسته \en{Coreutils}:

\begin{latin}
\url{http://www.gnu.org/software/coreutils/manual/coreutils.html#Output-of-entire-files}\\
\url{http://www.gnu.org/software/coreutils/manual/coreutils.html#Operating-on-sorted-files}\\
\url{http://www.gnu.org/software/coreutils/manual/coreutils.html#Operating-on-fields}\\
\url{http://www.gnu.org/software/coreutils/manual/coreutils.html#Operating-on-characters}
\end{latin}

از بسته \en{Diffutils}:

\begin{latin}
\url{http://www.gnu.org/software/diffutils/manual/html_mono/diff.html}
\end{latin}

ابزار \cmd{sed}:

\begin{latin}
\url{http://www.gnu.org/software/sed/manual/sed.htm}
\end{latin}

ابزار \en{Aspell}:

\begin{latin}
\url{http://aspell.net/man-html/index.html}
\end{latin}

منابع آنلاین متعدد دیگری نیز به طور خاص برای \cmd{sed} وجود دارد:

\begin{latin}
\url{http://www.grymoire.com/Unix/Sed.html}\\
\url{http://sed.sourceforge.net/sed1line.txt}
\end{latin}

علاوه بر این، جستجوی عباراتی نظیر \en{``sed one liners''} و \en{``sed cheat sheets''} در موتورهای جستجو، الگوهای عملیاتی بسیار ارزشمندی را در اختیار شما قرار می‌دهد.

\section{تمرین‌های تکمیلی}

چند ابزار کلیدی دیگر در حوزه‌ی پردازش متن وجود دارند که بررسی آن‌ها برای تسلط بیشتر توصیه می‌شود:

\begin{itemize}
  \item \cmd{split}: تقسیم یک فایل حجیم به قطعات کوچک‌تر بر اساس تعداد سطور یا حجم فایل.
  \item \cmd{csplit}: تفکیک فایل به بخش‌های مختلف بر اساس الگوهای محتوایی (\en{Context}).
  \item \cmd{sdiff}: مقایسه‌ی دو فایل و نمایش تفاوت‌ها به‌صورت ستون‌های کنار هم جهت ادغام بصری و دقیق‌تر.
\end{itemize}